\SourcePage{163}
\SGAsection{6}{不必仿射的代数群的极大环面与 Cartan 子群（代数闭基域）}
\begin{SGAresult}{\Term{lemma}{引理}}{6.1}
设 $G$ 是域 $k$ 上的\Term{connected}{连通}\Term{algebraic-group}{代数群}，$Z$ 是其\Term{center}{中心}的\Term{finite-index}{有限指数}\Term{algebraic-subgroup}{代数子群}，则 $G/Z$ 是\Term{affine}{仿射的}。
\end{SGAresult}
\begin{proof}
可以假定 $Z$ 就是 $G$ 的\Term{center}{中心}，因为\Term{affine}{仿射}\Term{scheme}{概形}上的\Term{finite}{有限}\Term{scheme}{概形}仍是\Term{affine}{仿射的}。考虑\Term{vector-space}{向量空间}
\[
P^n=P^n(G)=\mathcal O_{G,e}/\mathfrak m_{G,e}^{n+1}\qquad(n\geq0\text{ 为整数}),
\]
其中 $\mathfrak m_{G,e}$ 为\Term{maximal-ideal}{极大理想}。于是 $G$ 通过\Term{adjoint-representation}{伴随表示}\Term{action}{作用}于 $P^n(G)$。若 $Z_n$ 是相应\Term{homomorphism}{同态}
\[
\operatorname{ad}_n:G\longrightarrow\operatorname{GL}(P^n)
\]
的\Term{kernel}{核}，利用 $G$ 的\Term{connected}{连通性}容易验证，$Z$ 是所有 $Z_n$ 的交，因而（$G$ 是\Term{noetherian}{诺特的}）等于其中某个 $Z_n$。但 $\operatorname{ad}_n$ 取商后给出\Term{monomorphism}{单态射} $G/Z_n\to\operatorname{GL}(P^n)$，故它是\Term{closed-immersion}{闭浸入}，从而 $G/Z_n$ 是\Term{affine}{仿射的}，$G/Z$ 也就如此。
\end{proof}

\SourcePage{164}
\begin{SGAresult}{\Term{lemma}{引理}}{6.2}
设 $G$ 是域 $k$ 上的\Term{smooth}{光滑}\Term{algebraic-group}{代数群}，$Z$ 是一个\Term{central-subgroup}{中心代数子群}，$G'=G/Z$，$u:G\to G'$ 为\Term{canonical}{典范}\Term{homomorphism}{同态}，$T$ 是 $G$ 中的\Term{multiplicative-type}{乘法型子群}，$T'=u(T)$ 为像群，$C$ 是 $T$ 在 $G$ 中的\Term{centralizer}{中心化子}，$C'$ 是 $T'$ 在 $G'$ 中的\Term{centralizer}{中心化子}。则有
\[
C'{}^0\subset u(C).
\]
\end{SGAresult}
\begin{proof}
可以假定 $k$ 是\Term{algebraically-closed}{代数闭}的。令
\[
C_1=(u^{-1}(C')^0)_{\mathrm{red}}.
\]
只需证明 $C_1\subset C$：因为 $u(C_1)$ 是 $C'$ 中的\Term{connected}{连通}\Term{finite-index}{有限指数}\Term{subgroup}{子群}，故在\Term{set}{集合}意义下等于 $C'{}^0$；又因 $C'$，从而 $C'{}^0$ 都是\Term{smooth}{光滑的}，二者便相等。考虑由 $G\times G$ 到 $G$ 的\Term{morphism}{态射}
\[
(g,t)\longmapsto gtg^{-1}t^{-1}.
\]
它诱导\Term{morphism}{态射}
\[
\varphi:C_1\times T\longrightarrow Z_1,\qquad Z_1=Z_{\mathrm{red}}.
\]
事实上，因左端是\Term{reduced}{约化的}，只需证明对 $g\in C_1(k)$、$t\in T(k)$，有 $gtg^{-1}t^{-1}\in Z(k)$；这来自 $g$ 模 $Z$ \Term{centralizes}{中心化} $T$。在 $k$ 上的\Term{rational-point}{有理点}处计算，容易看出 $\varphi(g,t)$ 对 $g$、对 $t$ 都是\Term{additive}{加性}的，因而是“\Term{bilinear}{双线性}的”。我们断言，由于 $Z_1$ \Term{smooth}{光滑}而 $C_1$ \Term{connected}{连通}，这一\Term{homomorphism}{同态}恒为零；这就证明 $C_1\subset C$。利用 $T$ 的\Term{dense}{稠密}性\Term{theorem}{定理}，可归结到 $T$ \Term{finite}{有限}的情形，即存在整数 $n>0$ 使 $n\operatorname{id}_T=0$ 的情形。注意 $\varphi$ 由一个群\Term{homomorphism}{同态}
\[
C_1\longrightarrow\operatorname{Hom}_{S\text{-gr}}(T,Z_1)
\]
给出。由于 $Z_1$ 交换且\Term{smooth}{光滑}，右端由 $k$ 上的\Term{etale}{平展}\Term{algebraic-group}{代数群}表示；而 $C_1$ \Term{connected}{连通}，故从 $C_1$ 到该群的任意群\Term{homomorphism}{同态}均为零。
\end{proof}
\begin{SGAresult}{推论}{6.3}
在上述条件下，假定 $C'$ \Term{connected}{连通}，则
\[
C'=u(C),\qquad C=u^{-1}(C').
\]
\end{SGAresult}
\begin{proof}
这时 $C'=C'{}^0$；另一方面，$C$ 显然包含 $Z$，故等于 $u^{-1}(u(C))$。
\end{proof}
\begin{SGAresult}{\Term{lemma}{引理}}{6.4}
设 $G$ 是域 $k$ 上的\Term{algebraic-group}{代数群}，$Z$ 是\Term{central-subgroup}{中心代数子群}，使得 $G/Z=G'$ 是\Term{torus}{环面}。则 $G$ 交换；而且若 $k$ \Term{algebraically-closed}{代数闭}，则 $G$ 中存在\Term{torus}{环面} $T$ 使 $u(T)=G'$，其中 $u:G\to G/Z=G'$ 为\Term{canonical}{典范}\Term{homomorphism}{同态}。
\end{SGAresult}
\begin{proof}
可以假定 $k$ \Term{algebraically-closed}{代数闭}。再次考虑由 $(g,h)\mapsto ghg^{-1}h^{-1}$ 给出的\Term{morphism}{态射} $G\times G\to G$。因 $Z$ 在\Term{center}{中心}中，该\Term{morphism}{态射}通过 $G'\times G'\to G$ \Term{resolution}{分解}；因 $G'$ 交换，后者取值于 $Z$，甚至取值于 $Z_1=Z_{\mathrm{red}}$，因为 $G'\times G'$ 是\Term{reduced}{约化的}。同上可见，所得\Term{morphism}{态射} $\varphi:G'\times G'\to Z_{\mathrm{red}}$ 为\Term{bilinear}{双线性}，故为零，这证明 $G$ 交换。为寻找 $G$ 中满足 $u(T)=G'$ 的\Term{torus}{环面} $T$，可以将 $G$ 换成 $G^0_{\mathrm{red}}$，从而假定 $G$ \Term{smooth}{光滑}且\Term{connected}{连通}；
\SourcePage{165}
还可将 $Z$ 换成 $Z^0_{\mathrm{red}}$，从而假定 $Z$ \Term{smooth}{光滑}且\Term{connected}{连通}。根据 Chevalley 的一个著名\Term{theorem}{定理}\footnote{编者注：此处应给出参考文献。\SGAsourcegap{原书此处未给出所引用 Chevalley \Term{theorem}{定理}的参考文献。}}，$Z$ 是\Term{abelian-variety}{阿贝尔簇}被\Term{affine}{仿射}群的\Term{extension}{扩张}，于是立即归结为在下述两种情形分别证明断言：$1^\circ)$ $Z$ \Term{affine}{仿射}；$2^\circ)$ $Z$ 是\Term{abelian-variety}{阿贝尔簇}。在情形 $1^\circ)$，$G$ 是\Term{affine}{仿射的}，结论众所周知（BIBLE 7 \Term{theorem}{定理} 3 a)）。因此假定 $Z$ 为\Term{abelian-variety}{阿贝尔簇}。从\Term{additive-group}{加法群} $\mathbb G_a$ 到\Term{torus}{环面} $G'=R$ 或到\Term{abelian-variety}{阿贝尔簇} $Z$ 的任意\Term{homomorphism}{同态}均为平凡的，故从 $\mathbb G_a$ 到 $G$ 的任意\Term{homomorphism}{同态}亦然，因此 $G$ 不含\Term{isomorphism}{同构}于 $\mathbb G_a$ 的\Term{subgroup}{子群}。由已经用过的 Chevalley 结构\Term{theorem}{定理}，有\Term{exact-sequence}{正合列}
\[
0\longrightarrow T\longrightarrow G\longrightarrow A\longrightarrow0,\tag{*}
\]
其中 $A$ 是\Term{abelian-variety}{阿贝尔簇}，$T$ 是\Term{smooth}{光滑}\Term{connected}{连通}\Term{affine}{仿射}群。后者交换且不含\Term{additive-group}{加法子群}，故 $T$ 是\Term{torus}{环面}（而且显然是 $G$ 的唯一\Term{maximal-torus}{极大环面}）。问题归结为证明，每个\Term{epimorphism}{满同态}
\[
u:G\longrightarrow R,
\]
其中 $R$ 为\Term{torus}{环面}，都满足 $u(T)=R$。令
\[
\operatorname{Hom}_{\mathrm{gr}}(T,\mathbb G_m)=M,\qquad
\operatorname{Hom}_{\mathrm{gr}}(R,\mathbb G_m)=P
\]
（它们是\Term{finitely-generated}{有限生成}自由 $\mathbb Z$-模，并由 $T=D_k(M)$、$R=D_k(P)$ 在\Term{isomorphism}{同构}意义下恢复 $T,R$），又令
\[
B=\operatorname{Ext}^1_{k\text{-gr}}(A,\mathbb G_m)
\]
（$B$ 也就是 $A$ 的\Term{duality}{对偶}\Term{abelian-variety}{阿贝尔簇}的 $k$-\Term{rational-point}{有理点}集）。显然有
\[
\operatorname{Ext}^1(A,T)=\operatorname{Hom}_{\mathrm{gr}}(M,B),\qquad
\operatorname{Ext}^1(A,R)=\operatorname{Hom}_{\mathrm{gr}}(P,B).\tag{xx}
\]
特别地，$A$ 被 $T$ 的\Term{extension}{扩张} $G$ 由一个\Term{homomorphism}{同态}
\[
\theta:M\longrightarrow B
\]
给出。另一方面，\Term{exact-sequence}{正合列} $(*)$ 给出\Term{exact-sequence}{正合列}
\[
0\longrightarrow\operatorname{Hom}(A,R)\longrightarrow\operatorname{Hom}(G,R)
\longrightarrow\operatorname{Hom}(T,R)\longrightarrow\operatorname{Ext}^1(A,R).
\]
此外 $\operatorname{Hom}(A,R)=0$，而 $\operatorname{Hom}(T,R)\to\operatorname{Ext}^1(A,R)$ 等同于由 $\theta:M\to B$ \Term{derive}{导出}的\Term{homomorphism}{同态}
\[
\operatorname{Hom}(P,M)\longrightarrow\operatorname{Hom}(P,B).
\]
令 $N=\operatorname{Ker}(\theta)$，便得到\Term{canonical}{典范}\Term{bijection}{双射}
\[
\operatorname{Hom}(G,R)\simeq\operatorname{Hom}(P,N)\simeq\operatorname{Hom}(S,R),
\qquad S=D_k(N).
\]
其直接的几何描述如下。
\end{proof}
\SourcePage{166}
\begin{SGAresult}{\Term{lemma}{引理}}{6.5}
设 $G$ 是\Term{abelian-variety}{阿贝尔簇} $A$ 被\Term{torus}{环面} $T=D_k(M)$ 的\Term{extension}{扩张}，由\Term{homomorphism}{同态} $\theta:M\to B=\operatorname{Ext}^1(A,\mathbb G_m)$ 定义（\Term{base-field}{基域} $k$ \Term{algebraically-closed}{代数闭}）。令 $N=\operatorname{Ker}\theta$，$S=D_k(N)$ 为相应\Term{torus}{环面}，它\Term{isomorphism}{同构}于 $T/U$，其中 $U=D_k(M/N)$；考虑 $A$ 被 $S$ 的\Term{extension}{扩张} $H=G/U$。该\Term{extension}{扩张}是\Term{splitting}{分裂}的\footnote{编者注：原文写作 ``splitte''，应为 ``est scindée''（\Term{splitting}{分裂}）。}，因此有唯一的从 $H$ 到 $S$ 的\Term{projection}{投影}，进而有唯一的\Term{homomorphism}{同态}
\[
v:G\longrightarrow S
\]
延拓\Term{canonical}{典范}\Term{homomorphism}{同态} $T\to S$。于是，对任意\Term{torus}{环面} $R$ 及任意\Term{homomorphism}{同态} $u:G\to R$，存在唯一的\Term{homomorphism}{同态} $u':S\to R$ 使 $u=u'v$。特别地，$\operatorname{Im}(u)=\operatorname{Im}(u')=u(T)$。
\end{SGAresult}
这特别说明，若 $u$ 是\Term{epimorphism}{满同态}，则其在 $T$ 上的\Term{restriction}{限制}亦然，从而完成了 6.4 的证明。
\begin{SGAresult}{\Term{theorem}{定理}}{6.6}
设 $G$ 是\Term{algebraically-closed-field}{代数闭域} $k$ 上的\Term{smooth}{光滑}\Term{connected}{连通}\Term{algebraic-group}{代数群}。

 a) $G$ 的\Term{maximal-torus}{极大环面}彼此\Term{conjugacy}{共轭}。

 b) 若 $T$ 是 $G$ 的\Term{torus}{环面}，则其\Term{centralizer}{中心化子}是\Term{smooth}{光滑}\Term{connected}{连通的}。

 c) 映射 $T\mapsto\operatorname{Centr}_G(T)$ 给出 $G$ 的\Term{maximal-torus}{极大环面} $T$ 的\Term{set}{集合}与 $G$ 的\Term{cartan-subgroup}{Cartan 子群}（第 1 节）的\Term{set}{集合}之间的一一对应。$G$ 的\Term{algebraic-subgroup}{代数子群} $C$ 为\Term{cartan-subgroup}{Cartan 子群}，当且仅当它\Term{smooth}{光滑}、\Term{nilpotent}{幂零}，且在\Term{set}{集合}意义下等于其\Term{connected-normalizer}{连通正规化子}（这时它甚至等于其\Term{connected-normalizer}{连通正规化子}）；此时有 $C=\operatorname{Centr}_G(T)$，其中 $T$ 是 $C$ 的唯一\Term{maximal-torus}{极大环面}，并且 $\operatorname{Norm}_G(C)=\operatorname{Norm}_G(T)$。

 d) 设 $v:G\to H$ 为\Term{smooth}{光滑}\Term{connected}{连通}\Term{algebraic-group}{代数群}的\Term{epimorphism}{满同态}，则 $H$ 的\Term{maximal-torus}{极大环面}（相应地，\Term{cartan-subgroup}{Cartan 子群}）恰为 $G$ 的\Term{maximal-torus}{极大环面}（相应地，\Term{cartan-subgroup}{Cartan 子群}）的像。若 $T$ 是 $G$ 的\Term{maximal-torus}{极大环面}，$C$ 是其\Term{centralizer}{中心化子}，则 $v(C)$ 为 $v(T)$ 的\Term{centralizer}{中心化子}。

 e) 在 d) 的条件下，再假定 $\operatorname{Ker}v$ 是 $G$ 的\Term{central-subgroup}{中心子群}，则 $T\mapsto v(T)$（相应地，$C\mapsto v(C)$）给出 $G$ 的\Term{maximal-torus}{极大环面}（相应地，\Term{cartan-subgroup}{Cartan 子群}）与 $H$ 的\Term{maximal-torus}{极大环面}（相应地，\Term{cartan-subgroup}{Cartan 子群}）之间的一一对应。$G$ 的\Term{cartan-subgroup}{Cartan 子群}包含 $G$ 的\Term{center}{中心}，因而也包含 $\operatorname{Ker}v$，且恰为形如 $v^{-1}(C')$ 的群，其中 $C'$ 为 $H$ 的\Term{cartan-subgroup}{Cartan 子群}。
\end{SGAresult}
也立即写出 c) 的以下直接推论。
\begin{SGAresult}{推论}{6.7}
$G$ \Term{nilpotent}{幂零}（即群 $G(k)$ \Term{nilpotent}{幂零}），当且仅当 $G$ 中的\Term{torus}{环面}均在\Term{center}{中心}中；也\Term{equivalence}{等价}于 $G$ 仅有一个\Term{maximal-torus}{极大环面}（这时它就是 $G$ 的最大\Term{subtorus}{子环面}）。
\end{SGAresult}
\begin{proof}[6.6 的证明]
设 $Z$ 是 $G$ 的\Term{central-subgroup}{中心代数子群}，$G'=G/Z$，$u:G\to G'$ 为\Term{canonical}{典范}\Term{homomorphism}{同态}。则 $G'$ 是\Term{smooth}{光滑}\Term{connected}{连通}群。若 $T'$ 为 $G'$ 的\Term{subtorus}{子环面}，由 6.4，$u^{-1}(T')$ 交换，
\SourcePage{167}
且 $T'$ 为 $u^{-1}(T')$ 的某个\Term{subtorus}{子环面}、因而为 $G$ 的某个\Term{subtorus}{子环面}的像。由于 $u^{-1}(T')$ 交换，它显然有一个最大\Term{subtorus}{子环面} $T$（两个\Term{subtorus}{子环面}之和给出包含二者的第三个\Term{subtorus}{子环面}），并有 $u(T)=T'$。这立即说明，$G$ 的任意\Term{maximal-torus}{极大环面} $T$ 的像 $T'=u(T)$ 是 $G'$ 的\Term{maximal-torus}{极大环面}，且 $T\mapsto u(T)$ 为两者\Term{maximal-torus}{极大环面}间的一一对应。

现在取 $Z=\operatorname{Centr}(G)^0_{\mathrm{red}}$；由 6.1，$G'$ 是\Term{affine}{仿射的}。因此 $G'$ 的\Term{maximal-torus}{极大环面}彼此\Term{conjugacy}{共轭}，$G$ 的\Term{maximal-torus}{极大环面}亦然，a) 得证。此外，$G$ \Term{nilpotent}{幂零}，相应地仅有一个\Term{maximal-torus}{极大环面}，当且仅当 $G'$ 满足同一条件；$G'$ \Term{affine}{仿射}，故对 $G'$ 这两个条件\Term{equivalence}{等价}（BIBLE 6 \Term{theorem}{定理} 4 推论 2），对 $G$ 亦然。又若 $G$ 仅有一个\Term{maximal-torus}{极大环面} $T$，则它是\Term{normal-subgroup}{正规子群}，因而在\Term{center}{中心}中；$G$ 中任意\Term{torus}{环面}都包含在一个\Term{maximal-torus}{极大环面}中，故也在\Term{center}{中心}中。反之，若所有\Term{torus}{环面}都在\Term{center}{中心}中，则\Term{maximal-torus}{极大环面}亦然，由\Term{conjugacy}{共轭}\Term{theorem}{定理} a)，\Term{maximal-torus}{极大环面}便只有一个。这证明了 6.7。

设 $T$ 是 $G$ 的任意\Term{torus}{环面}，$T'=u(T)$，则 $C'=\operatorname{Centr}_{G'}(T')$ \Term{connected}{连通}（BIBLE 6 \Term{theorem}{定理} 6 a)）。因此由 6.3，$T$ 的\Term{centralizer}{中心化子} $C$ 等于 $u^{-1}(C)$，故\Term{connected}{连通}（因为 $Z$ \Term{connected}{连通}），b) 得证。若 $T$ 极大，因而 $T'$ 极大，则已知 $C'$ \Term{nilpotent}{幂零}，所以 $C$ 作为 $C'$ 的\Term{central-extension}{中心扩张}也\Term{nilpotent}{幂零}。又 $T$ 是 $C$ 的\Term{maximal-torus}{极大环面}，由 6.7，它是 $C$ 的唯一\Term{maximal-torus}{极大环面}。故映射 $T\mapsto\operatorname{Centr}_G(T)$ 从 $G$ 的\Term{maximal-torus}{极大环面}\Term{set}{集合}到\Term{cartan-subgroup}{Cartan 子群}\Term{set}{集合}是\Term{bijection}{双射}。此外有
\[
\operatorname{Centr}_G(T)=C\subset\operatorname{Norm}_G(C)\subset\operatorname{Norm}_G(T).
\]
已知 $T$ 的\Term{centralizer}{中心化子}在其\Term{normalizer}{正规化子}中具有\Term{finite-index}{有限指数}（参见 XI 5.9；该处论证不需\Term{affine}{仿射}假设，只需所讨论的两个\Term{functor}{函子}\Term{representable-functor}{可表示}，正如 XI 6.5 所指出），且 $C$ \Term{smooth}{光滑}\Term{connected}{连通}，故
\[
C=\operatorname{Norm}_G(C)^0.
\]
再由\Term{maximal-torus}{极大环面}与\Term{cartan-subgroup}{Cartan 子群}间的一一对应，可见 $\operatorname{Norm}_G(C)$ 和 $\operatorname{Norm}_G(T)$ 有相同的 $k$-值点；因后者\Term{smooth}{光滑}，故
\[
\operatorname{Norm}_G(T)=\operatorname{Norm}_G(C).
\]
为完成 c)，只需证明：若 $C$ 是 $G$ 的\Term{smooth}{光滑}\Term{nilpotent}{幂零}\Term{connected}{连通}\Term{subgroup}{子群}，在其\Term{normalizer}{正规化子}中具有\Term{finite-index}{有限指数}，则 $C$ 是\Term{cartan-subgroup}{Cartan 子群}。由于 $Z$ 在\Term{center}{中心}中，$C$ 的\Term{normalizer}{正规化子}包含 $Z$；而 $Z$ \Term{smooth}{光滑}\Term{connected}{连通}，故 $Z\subset C$，于是 $C=u^{-1}(C')$，其中 $C'=u(C)$。这时
\[
\operatorname{Norm}_G(C)=u^{-1}(\operatorname{Norm}_{G'}(C')),
\]
\SourcePage{168}
这证明 $C'$ \Term{nilpotent}{幂零}\Term{connected}{连通}，且在其\Term{normalizer}{正规化子}中具有\Term{finite-index}{有限指数}；故由 BIBLE 7 \Term{theorem}{定理} 1，$C'$ 是 $G'$ 的\Term{cartan-subgroup}{Cartan 子群}，立即得 $C$ 是 $G$ 的\Term{cartan-subgroup}{Cartan 子群}。

证明 e)。我们处于证明开头的条件下（令 $\operatorname{Ker}u=Z$、$G'=H$），已知 $T\mapsto u(T)$ 给出 $G$ 与 $G'$ 的\Term{maximal-torus}{极大环面}间的一一对应。由刚证明的\Term{maximal-torus}{极大环面}与\Term{cartan-subgroup}{Cartan 子群}间的一一对应，得到 $G$ 的\Term{cartan-subgroup}{Cartan 子群} $C$ 与 $G'$ 的\Term{cartan-subgroup}{Cartan 子群} $C'$ 间的一一对应，其中对应于 $C=\operatorname{Centr}_G(T)$ 的是 $C'=\operatorname{Centr}_{G'}(T')$，$T'=u(T)$。由 b)，$C'$ \Term{connected}{连通}，故由 6.3，$C'=u(C)$ 且 $C=u^{-1}(C')$，e) 得证。

最后证明 d)。设 $T$ 是 $G$ 的\Term{maximal-torus}{极大环面}，证明 $v(T)$ 是 $H$ 的\Term{maximal-torus}{极大环面}（由\Term{conjugacy}{共轭}\Term{theorem}{定理} a)，这就说明 $H$ 的\Term{maximal-torus}{极大环面}全都形如这样的 $v(T)$）。设 $R$ 是 $H$ 中包含 $v(T)$ 的\Term{torus}{环面}，要证 $R=v(T)$。将 $H$ 换为 $R$、$G$ 换为 $v^{-1}(R)^0_{\mathrm{red}}$，可假定 $R=H$，即 $H$ 是\Term{torus}{环面}，问题归结为证明 $v(T)=H$。仍令 $Z=\operatorname{Centr}(G)^0_{\mathrm{red}}$、$G'=G/Z$、$H'=H/v(Z)$，得到图
\[
\begin{tikzcd}
e\arrow[r]&Z\arrow[r]\arrow[d,"v''"]&G\arrow[r,"u"]\arrow[d,"v"]&G'\arrow[r]\arrow[d,"v'"]&e\\
e\arrow[r]&v(Z)\arrow[r]&H\arrow[r,"u'"]&H'\arrow[r]&e.
\end{tikzcd}\tag{D}
\]
已知 $u(T)$ 是 $G'$ 的\Term{maximal-torus}{极大环面}；而 $G'$ \Term{affine}{仿射}，故 $v':G'\to H'$ 是\Term{smooth}{光滑}\Term{connected}{连通}\Term{affine}{仿射}群的\Term{epimorphism}{满同态}，所以 $v'(u(T))$ 是 $H'$ 的\Term{maximal-torus}{极大环面}（BIBLE 7 \Term{theorem}{定理} 3 a)），从而等于 $H'$。即 $u'(v(T))=H'$，故为证 $v(T)=H$，只需证明 $v(T)\supset v(Z)$。但 $v(Z)$ 是 $H$ 的\Term{smooth}{光滑}\Term{connected}{连通}\Term{subgroup}{子群}，因而是\Term{torus}{环面}；$v'':Z\to v(Z)$ 为\Term{epimorphism}{满同态}，故由 6.4，$v(Z)=v''(S)$，其中 $S$ 是 $Z$ 的某个\Term{torus}{环面}。$S$ 在 $G$ 的\Term{center}{中心}中，显然包含在\Term{maximal-torus}{极大环面} $T$ 中，故 $v(Z)\subset v(T)$。这证明了 d) 关于\Term{maximal-torus}{极大环面}的断言。

由\Term{maximal-torus}{极大环面}与\Term{cartan-subgroup}{Cartan 子群}之间的一一对应，只需证明：若 $T$ 是 $G$ 的\Term{maximal-torus}{极大环面}，$C$ 是其\Term{centralizer}{中心化子}，则 $v(C)$ 为 $v(T)$ 的\Term{centralizer}{中心化子}。为此，再用上图 (D)（此时当然不再假定 $H$ 是\Term{torus}{环面}），令 $T'=u(T)$、$C'=u(C)$。在 e) 中已知 $C'$ 是\Term{maximal-torus}{极大环面} $T'$ 的\Term{centralizer}{中心化子}；由于 $G'$ \Term{affine}{仿射}，因而 $H'$ 也\Term{affine}{仿射}，故 $v'(C')$ 是 $H'$ 的\Term{maximal-torus}{极大环面} $v'(T')$ 的\Term{centralizer}{中心化子}（BIBLE 7 \Term{theorem}{定理} 3 a)）。即 $u'(v(C))$ 是 $u'(v(T))$ 的\Term{centralizer}{中心化子}。由于 $C$ 包含 $Z$，故 $v(C)$ 包含 $v(Z)$，于是 $v(C)$ 是 $u'(v(C))$，即 $u'(v(T))$ 的\Term{centralizer}{中心化子}在 $u'$ 下的\Term{inverse-image}{逆像}；将 e) 应用于 $u':H\to H'$ 及 $H$ 的\Term{maximal-torus}{极大环面} $v(T)$，便知它是 $v(T)$ 的\Term{centralizer}{中心化子}。6.6 证毕。
\end{proof}

\SourcePage{169}
\SGAsection{7}{对不必仿射的光滑群预概形的应用}
\begin{SGAresult}{\Term{theorem}{定理}}{7.1}
设 $S$ 是\Term{prescheme}{预概形}，$G$ 是 $S$ 上\Term{smooth}{光滑}、\Term{separated}{分离}、\Term{finite-type}{有限型}的\Term{group-prescheme}{群预概形}。假定 $G$ 对于\Term{fpqc-topology}{忠实平坦拟紧拓扑}\Term{local-scheme}{局部}地有\Term{maximal-torus}{极大环面}。则：

 a) 映射 $T\mapsto\operatorname{Centr}_G(T)$ 诱导 $G$ 的\Term{maximal-torus}{极大环面}\Term{set}{集合}与 $G$ 的\Term{cartan-subgroup}{Cartan 子群}\Term{set}{集合}之间的\Term{bijection}{双射}。若 $C$ 对应于 $T$，则 $T$ 为 $C$ 的唯一\Term{maximal-torus}{极大环面}。

 b) 设 $T,T'$ 是 $G$ 的两个\Term{maximal-torus}{极大环面}，$C,C'$ 为相应\Term{cartan-subgroup}{Cartan 子群}，则
\[
\operatorname{Transp}_G(T,T')=\operatorname{Transp}_G(C,C')=\operatorname{Transp}_G(T,C').
\]
前两项也等于相应的\Term{strict-transporter}{严格输送子}；这一\Term{functor}{函子}由 $G$ 的一个在 $S$ 上\Term{smooth}{光滑}的\Term{closed-subprescheme}{闭子预概形}表示。\Term{torus}{环面} $T,T'$ 以及\Term{cartan-subgroup}{Cartan 子群} $C,C'$ 对于\Term{etale-topology}{平展拓扑}\Term{local-scheme}{局部}\Term{conjugacy}{共轭}。

 c) 对于\Term{etale-topology}{平展拓扑}，\Term{local-scheme}{局部}地存在 $G$ 的\Term{maximal-torus}{极大环面}和\Term{cartan-subgroup}{Cartan 子群}。

 d) 假定 $G$ 的任一\Term{fiber}{纤维} $G_s$ 的每个\Term{finite}{有限}\Term{subset}{子集}都包含在 $G$ 的一个\Term{affine-open}{仿射开集}内（例如 $G$ 在 $S$ 上\Term{quasi-projective}{拟射影}，或 $S$ 为\Term{artinian}{阿廷的}），则 1.10 中定义的\Term{functor}{函子} $\mathcal T:(\mathrm{Sch})^\circ_{/S}\to(\mathrm{Ens})$（$G$ 的\Term{maximal-torus}{极大环面函子}），它\Term{isomorphism}{同构}于 $G$ 的\Term{cartan-subgroup}{Cartan 子群}\Term{functor}{函子}，由 $S$ 上\Term{smooth}{光滑}、\Term{separated}{分离}、\Term{finite-type}{有限型}的\Term{prescheme}{预概形}表示；当 $G$ 在 $S$ 上\Term{quasi-projective}{拟射影}时，它也\Term{quasi-projective}{拟射影}；当 $G$ 在 $S$ 上\Term{affine}{仿射}或 $S$ 为\Term{artinian}{阿廷的}时，它在 $S$ 上\Term{affine}{仿射}。

 e) 设 $u:G\to G'$ 是 $S$-\Term{group-prescheme}{群预概形}的\Term{morphism}{态射}，其中 $G'$ 在 $S$ 上\Term{smooth}{光滑}、\Term{separated}{分离}且\Term{finite-type}{有限型}；假定对任意 $s\in S$，有 $u_s(G_s)=G'_s$，即 $u_s$ \Term{faithfully-flat}{忠实平坦}（第 VI 讲\footnote{编者注：未在第 VI$_{\mathrm B}$ 讲找到此引用，参见 M. Demazure 与 P. Gabriel，\textit{Groupes algébriques}，I，Masson（1970），命题 II.5.1.(c)。}）。则对 $G$ 的任意\Term{maximal-torus}{极大环面} $T$，$u(T)$ 是 $G'$ 的\Term{maximal-torus}{极大环面}；若 $G'$ 的\Term{fiber}{纤维}\Term{connected}{连通}，则对 $G$ 的任意\Term{cartan-subgroup}{Cartan 子群} $C$，$u(C)$ 为 $G'$ 的\Term{cartan-subgroup}{Cartan 子群}；若 $C$ 是 $T$ 的\Term{centralizer}{中心化子}，则 $u(C)$ 是 $u(T)$ 的\Term{centralizer}{中心化子}。这两种情形下，诱导的\Term{morphism}{态射} $T\to u(T)$，相应地 $C\to u(C)$，都是\Term{faithfully-flat}{忠实平坦的}。

 f) 在 e) 的条件下，再假定 $\operatorname{Ker}u$ 为 $G$ 的\Term{central-subgroup}{中心子群}。则 $T\mapsto u(T)$ 是 $G$ 的\Term{maximal-torus}{极大环面}\Term{set}{集合}到 $G'$ 的\Term{maximal-torus}{极大环面}\Term{set}{集合}的\Term{bijection}{双射}；若 $G'$ 的\Term{fiber}{纤维}\Term{connected}{连通}，则 $C\mapsto u(C)$ 同样是 $G$ 的\Term{cartan-subgroup}{Cartan 子群}\Term{set}{集合}到 $G'$ 的\Term{cartan-subgroup}{Cartan 子群}\Term{set}{集合}的\Term{bijection}{双射}；若 $C'=u(C)$，则 $C=u^{-1}(C')$。
\end{SGAresult}
\SourcePage{170}
\textbf{注记 7.2。} 我们将在第 XV 讲看到，去掉对 $G_s$ 的\Term{finite}{有限}\Term{subset}{子集}的\Term{restriction}{限制}，d) 的结论仍成立。在那里还将证明：只假定 $G$ 对于\Term{fpqc-topology}{忠实平坦拟紧拓扑}\Term{local-scheme}{局部}地有\Term{cartan-subgroup}{Cartan 子群}（但不一定有\Term{maximal-torus}{极大环面}），则 b)、c)、d)、e) 中仅涉及\Term{cartan-subgroup}{Cartan 子群}的结论仍成立。为此还要用到 $S$ 为\Term{artinian}{阿廷的}时的 7.1。此外，使用第 XV 讲的方法，c) 和 d)（$S$ 不为\Term{artinian}{阿廷的}时）的证明会大大简化。

\begin{proof}[7.1 的证明]
a) 仿照 3.2，可见对 $G$ 的任意\Term{maximal-torus}{极大环面} $T$，$C=\operatorname{Centr}_G(T)$ 确为 $G$ 的\Term{cartan-subgroup}{Cartan 子群}，且 $T$ 由 $C$ 唯一确定，为 $C$ 的唯一\Term{maximal-torus}{极大环面}。因此只需证明 $G$ 的每个\Term{cartan-subgroup}{Cartan 子群} $C$ 都由 $G$ 的某个\Term{maximal-torus}{极大环面}定义，亦即 $C$ 有一个在\Term{center}{中心}中的\Term{maximal-torus}{极大环面}。问题对 \Term{fpqc-topology}{fpqc 拓扑}仍是\Term{local-scheme}{局部}的，故可假定 $G$ 有\Term{maximal-torus}{极大环面} $T$。由 XI 6.2，$\operatorname{Transp}_G(T,C)$ 由 $G$ 的一个在 $S$ 上\Term{smooth}{光滑}的\Term{closed-subprescheme}{闭子预概形} $S'$ 表示；又由 6.6 c)，$S'\to S$ \Term{surjectivity}{满射}。因此作\Term{base-change}{基变换} $S'\to S$ 后，可假定存在 $G$ 的\Term{section}{截面} $g$ 使 $\operatorname{ad}(g)\cdot T\subset C$。这时 $\operatorname{ad}(g)\cdot T$ 是 $C$ 的\Term{maximal-torus}{极大环面}，由 6.6 c) 它逐\Term{fiber}{纤维}处于\Term{center}{中心}中，故由 IX 5.6 b) 它在\Term{center}{中心}中。

b) 设 $T,T'$ 是 $G$ 的两个\Term{maximal-torus}{极大环面}，$C,C'$ 是 $G$ 的相应\Term{cartan-subgroup}{Cartan 子群}。则下述条件\Term{equivalence}{等价}：
\[
(1)\ T\subset T'\quad (2)\ T=T'\quad (3)\ T\subset C'\quad
(4)\ C\subset C'\quad (5)\ C=C'.
\]
这由 a) 显然得出。任意\Term{base-change}{基变换}后仍用此结果，便得到 b) 中各\Term{transporter}{输送子}和\Term{strict-transporter}{严格输送子}间的等式。此外，在 a) 中已经指出 $\operatorname{Transp}_G(T,C')$ 由 $G$ 的一个在 $S$ 上\Term{smooth}{光滑}的\Term{closed-subprescheme}{闭子预概形}表示，且其\Term{structure-morphism}{结构态射}\Term{surjectivity}{满射}。因此由 Hensel \Term{lemma}{引理}，该\Term{prescheme}{预概形}在 $S$ 上对于\Term{etale-topology}{平展拓扑}\Term{local-scheme}{局部}有\Term{section}{截面}。故 $T,T'$ 以及 $C,C'$ 分别对于\Term{etale-topology}{平展拓扑}\Term{local-scheme}{局部}\Term{conjugacy}{共轭}，b) 得证。

c) 先假定 $S$ 为\Term{local}{局部}\Term{artinian}{阿廷的}。当 $G$ 有\Term{maximal-torus}{极大环面} $T$ 时，由 b) 中证明的\Term{conjugacy}{共轭}\Term{theorem}{定理}，$G$ 的\Term{maximal-torus}{极大环面}\Term{functor}{函子} $\mathcal T$ 由\Term{homogeneous-space}{齐性空间} $G/N$ 表示，其中 $N$ 为 $T$ 在 $G$ 中的\Term{normalizer}{正规化子}，由 b) 它\Term{smooth}{光滑}。又如 VI$_{\mathrm A}$ 3.2.1 所述，由于 $\mathcal T$ 是 $G$ 下的\Term{homogeneous-space}{齐性空间}，$\mathcal T$ 的任意\Term{finite}{有限}\Term{subset}{子集}都包含在某个\Term{affine-open}{仿射开集}内。一般情形下，$S$ 的\Term{residue-field}{剩余域} $k$ 有一个\Term{finite-extension}{有限扩张} $k'$，使 $G_{k'}$ 有\Term{maximal-torus}{极大环面}；$k'$ 可由 $k$ 通过\Term{finite}{有限}\Term{flat}{平坦}\Term{base-change}{基变换} $S'\to S$ 得到，而 $G_{k'}$ 的\Term{maximal-torus}{极大环面}\Term{liftable}{可提升}为 $G_{S'}$ 的\Term{maximal-torus}{极大环面}（XI 2.1 bis）。故\Term{functor}{函子} $\mathcal T_{S'}$ 由 $S'$ 上\Term{smooth}{光滑}、\Term{separated}{分离}、\Term{finite-type}{有限型}的\Term{prescheme}{预概形}表示，且其每个\Term{finite}{有限}\Term{subset}{子集}包含在一个\Term{affine-open}{仿射开集}内。因此 $\mathcal T_{S'}$ 上的自然\Term{descent-datum}{下降资料}\Term{effective}{有效}，$\mathcal T$ \Term{representable-functor}{可表示}；由\Term{descent}{下降}，$\mathcal T$ 在 $S$ 上\Term{smooth}{光滑}、\Term{separated}{分离}且\Term{finite-type}{有限型}。由\Term{smooth}{光滑性}及 Hensel \Term{lemma}{引理}，$\mathcal T$ 对于
\SourcePage{171}
\Term{etale-topology}{平展拓扑}\Term{local-scheme}{局部}地有\Term{section}{截面}。这证明了 $S$ 为\Term{artinian}{阿廷的}时的 c)、d)（注意：稍后将证明这时 $\mathcal T$ 实际上在 $S$ 上\Term{affine}{仿射}）。

现在设 $S$ 任意。c)、d) 对 $S$ 的 Zariski \Term{topology}{拓扑}是\Term{local-scheme}{局部}断言，故可假定存在素数 $\ell$ 与 $S$ 的各\Term{residue-characteristic}{剩余特征}互素，$S$ \Term{affine}{仿射}，且 $G$ 的\Term{fiber}{纤维}的\Term{reductive-rank}{约化秩}恒定为 $r$（由 fpqc \Term{local-scheme}{局部}存在\Term{maximal-torus}{极大环面}的假设，该秩显然\Term{locally-constant}{局部常值}）。设 $S'\to S$ 为\Term{faithfully-flat}{忠实平坦}\Term{quasi-compact}{拟紧}\Term{morphism}{态射}，$S'$ \Term{affine}{仿射}，且 $G_{S'}$ 有\Term{maximal-torus}{极大环面} $T'$。令 $C'$ 为其\Term{centralizer}{中心化子}。我们断言，存在 $\ell$ 的适当幂 $n$ 使得还有 $C=\operatorname{Centr}_G({}_nT)$。为此立即可归结到 $S'$ \Term{noetherian}{诺特}的情形，而这已在 XI 6.2 中证明。固定这样的 $n$，令
\[
M=(\mathbb Z/n\mathbb Z)^r,\qquad P=\operatorname{Hom}_{S\text{-gr}}(M_S,G).
\]
则 $P$ 显然\Term{representable-functor}{可表示}为 $({}_nG)^r$ 的一个\Term{finite-presentation}{有限表现}\Term{closed-subprescheme}{闭子预概形}，其中
\[
{}_nG=\operatorname{Hom}_{S\text{-gr}}((\mathbb Z/n\mathbb Z)_S,G)
\]
也是 $G$ 中 $n$ 次幂\Term{morphism}{态射}的\Term{kernel}{核}，因而\Term{representable-functor}{可表示}为 $G$ 的一个\Term{finite-presentation}{有限表现}\Term{closed-subprescheme}{闭子预概形}。又由 XI 2.1，$P$ 在 $S$ 上\Term{smooth}{光滑}。设 $s\in S$。由上面已证的结论，$\kappa(s)$ 有\Term{finite-separable-extension}{有限可分扩张} $k''$，使 $G_{k''}$ 中存在\Term{maximal-torus}{极大环面} $T''$；且因 ${}_nT''$ 在 $k''$ 上\Term{etale}{平展}（$n$ 与 $k''$ 的\Term{characteristic}{特征}互素），将 $k''$ 再换为一个\Term{finite-separable-extension}{有限可分扩张}，可假定 ${}_nT''$ \Term{isomorphism}{同构}于 $M_{k''}$。设 $S''\to S$ 是\Term{etale}{平展}\Term{morphism}{态射}，在 $s\in S$ 上方有点 $s''\in S''$，并诱导\Term{residue-extension}{剩余扩张} $\kappa(s'')\simeq k''$。于是 $P_{S''}\otimes\kappa(s'')$ 在 $\operatorname{Spec}(\kappa(s''))$ 上有\Term{section}{截面}；由 Hensel \Term{lemma}{引理}，将 $(S'',s'')$ 换为其上的\Term{etale}{平展} $(S''',s''')$ 后，可假定 $P_{S''}$ 在 $S''$ 上有\Term{section}{截面}，即有 $P(S'')$ 中的元素延拓给定\Term{section}{截面}。换言之，有\Term{homomorphism}{同态} $M_{S''}\to G_{S''}$，\Term{induction-isomorphism}{诱导同构} $M_{\kappa(s'')}\simeq{}_nT''$。由 IX 6.4（此处只需应用 \Term{nakayama-lemma}{Nakayama 引理}），此\Term{homomorphism}{同态}在 $s'$ 的某个\Term{open-neighborhood}{开邻域}上为\Term{closed-immersion}{闭浸入}；可假定该\Term{neighborhood}{邻域}即 $S''$。设其像为 $H''$，考虑它在 $G_{S''}$ 中的\Term{centralizer}{中心化子} $C''$；由 XI 6.2，这是 $G$ 的\Term{smooth}{光滑}\Term{closed-subprescheme}{闭子预概形}。再注意以下\Term{lemma}{引理}。
\end{proof}
\begin{SGAresult}{\Term{lemma}{引理}}{7.3}
在上述关于 $\ell,n,r$ 的条件下，对 $S$ 上的任意\Term{prescheme}{预概形} $S''$ 及 $G_{S''}$ 的任意\Term{maximal-torus}{极大环面} $T''$，有
\[
\operatorname{Centr}_{G_{S''}}(T'')=\operatorname{Centr}_{G_{S''}}({}_nT'').
\]
\end{SGAresult}
\begin{proof}
从 $S'\times_SS''$ 作\Term{faithfully-flat}{忠实平坦}\Term{descent}{下降}，可归结到 $G_{S''}$ 有一个满足上述关系的\Term{maximal-torus}{极大环面} $T''_1$ 的情形。但由 b)，$T''$ 与 $T''_1$ 对于\Term{etale-topology}{平展拓扑}\Term{local-scheme}{局部}\Term{conjugacy}{共轭}，故 $T''$ 也满足同一关系。
\end{proof}
将上述结果应用于 $\operatorname{Spec}(\kappa(s''))$ 而非 $S''$，可见\Term{fiber}{纤维} $C''_{s''}$ 是 $G$ 的\Term{cartan-subgroup}{Cartan 子群}。现在注意，$G$ 在 $S$ 上\Term{smooth}{光滑}，故各\Term{fiber}{纤维} $G_s$ 的\Term{identity-component}{单位连通分支}之并是 $G$ 的\Term{open-subset}{开集}（由 EGA IV 关于\Term{smooth}{光滑}\Term{morphism}{态射}的一般结果\footnote{编者注：应明确给出此引用。\SGAsourcegap{原书未给出所用 EGA IV 结果的具体编号。}}），
\SourcePage{172}
显然在 $G$ 的\Term{group-law}{群运算}下\Term{stable}{稳定}，赋予 $G$ 诱导的\Term{prescheme}{预概形}结构后便为 $G$ 的\Term{subgroup}{子群}。此外，$G^0$ 满足对 $G$ 的预备假设，而 $G$ 与 $G^0$ 的\Term{maximal-torus}{极大环面}显然一一对应。因此，为证 c)、d)，可假定 $G$ 的\Term{fiber}{纤维}\Term{connected}{连通}，以下即如此。这时 $G$ 的\Term{cartan-subgroup}{Cartan 子群}的\Term{fiber}{纤维}\Term{connected}{连通}（6.6 a)）。我们断言，在 $s''$ 的某个\Term{open-neighborhood}{开邻域}上，$C''{}^0$ 是 $G_{S''}$ 的\Term{cartan-subgroup}{Cartan 子群}（这便完成 c) 的证明）。这来自下述\Term{lemma}{引理}。
\begin{SGAresult}{\Term{lemma}{引理}}{7.4}
在 7.1 的条件下，假定 $G$ 的\Term{fiber}{纤维}\Term{connected}{连通}。设 $C$ 是 $G$ 的\Term{closed-subgroup-prescheme}{闭子群预概形}，在 $S$ 上\Term{smooth}{光滑}，$s\in S$ 使 $C_s$ 为 $G_s$ 的\Term{cartan-subgroup}{Cartan 子群}。则在 $s$ 的某个\Term{open-neighborhood}{开邻域}上，$C^0$ 是 $G$ 的\Term{cartan-subgroup}{Cartan 子群}。
\end{SGAresult}
\begin{proof}
容易归结到 $S$ \Term{local}{局部}且 $s$ 为其\Term{closed-point}{闭点}的情形，此时只需证明 $C$ 是 $G$ 的\Term{cartan-subgroup}{Cartan 子群}；再由\Term{flat}{平坦}\Term{descent}{下降}，可假定 $G$ 有\Term{maximal-torus}{极大环面} $T$。由 XI 6.2，$\operatorname{Transp}_G(T,C)$ 由 $G$ 的一个在 $S$ 上\Term{smooth}{光滑}的\Term{closed-subprescheme}{闭子预概形}表示。由对 $C_s$ 的假设及\Term{conjugacy}{共轭}\Term{theorem}{定理} 6.6 a)，该\Term{transporter}{输送子}在 $s$ 处的\Term{fiber}{纤维}\Term{nonempty}{非空}。由\Term{faithfully-flat}{忠实平坦}\Term{descent}{下降}，可归结到该\Term{transporter}{输送子}在 $S$ 上有\Term{section}{截面}的情形，即 $C$ 包含 $G$ 的一个\Term{maximal-torus}{极大环面} $T$。由 IX 5.6 a)，$T$ 在 $C^0$ 的\Term{center}{中心}中，故 $C^0\subset\operatorname{Centr}_G(T)$。这是 $S$ 上\Term{smooth}{光滑}\Term{group-scheme}{群概形}之间的包含，两者\Term{relative-dimension}{相对维数}相同（等于其在 $s$ 处共同\Term{fiber}{纤维}的\Term{dimension}{维数}），且\Term{fiber}{纤维}\Term{connected}{连通}，故必相等，证明完成。
\end{proof}
\begin{proof}[7.1 的证明（续）]
d) 保留前面关于 $\ell,n,r$ 及 $G$ 的\Term{fiber}{纤维}\Term{connected}{连通}的记号与假设。设 $Q:(\mathrm{Sch})^\circ_{/S}\to(\mathrm{Ens})$ 是由下式定义的\Term{functor}{函子}：
\[
Q(S')=\{G_{S'}\text{ 中类型为 }(\mathbb Z/n\mathbb Z)^r=M\text{ 的乘法型子群}\}
\quad\text{（IX 1.4）}.
\]
则 $T\mapsto{}_nT$ 是一个\Term{morphism}{态射} $\varphi:\mathcal T\to Q$，由 7.3 它是\Term{monomorphism}{单态射}。断言 $Q$ 由 $S$ 上\Term{separated}{分离}\Term{finite-presentation}{有限表现}的\Term{prescheme}{预概形}表示。事实上，如 XI 3.12 a) 的证明中所述，有\Term{isomorphism}{同构}
\[
Q\simeq P'/\Gamma,
\]
其中 $P'$ 是 c) 中引入的\Term{prescheme}{预概形} $P=\operatorname{Hom}_{S\text{-gr}}(M_S,G)$ 的\Term{open-and-closed}{开闭}\Term{subprescheme}{子预概形}，对应于\Term{monomorphism}{单态射} $M_{S'}\to G_{S'}$（参见 IX 6.8），且 $\Gamma=(\operatorname{Aut}_{\mathrm{gr}}(M))_S$。假设 $G_s$ 的每个\Term{finite}{有限}\Term{subset}{子集}都包含在 $G$ 的某个\Term{affine-open}{仿射开集}内，这一性质对\Term{closed-subprescheme}{闭子预概形}及\Term{cartesian-product}{笛卡尔积}\Term{stable}{稳定}，故显然传给 $({}_nG)^r$，进而传给 $P$ 与 $P'$。
\SourcePage{173}
因此 $Q$ 由 $S$ 上的\Term{finite-presentation}{有限表现}\Term{prescheme}{预概形}表示（参见第 V 讲）\footnote{编者注：应明确给出此引用。\SGAsourcegap{原书未给出第 V 讲中所用商\Term{representability}{可表示性}结果的具体编号。}}。同理可见，若 $G$ 在 $S$ 上\Term{quasi-projective}{拟射影}（相应地，\Term{affine}{仿射}），则 $Q$ 亦然。无论如何，$Q$ 在 $S$ 上\Term{separated}{分离}。现在有下述\Term{lemma}{引理}。
\end{proof}
\begin{SGAresult}{\Term{lemma}{引理}}{7.5}
上述\Term{homomorphism}{同态} $\varphi=\mathcal T\to Q$ \Term{representable-functor}{可表示}为\Term{open-immersion}{开浸入}。
\end{SGAresult}
\begin{proof}
换言之，需证明：若 $H$ 为 $G$ 中类型为 $M$ 的\Term{multiplicative-type}{乘法型子群}，则存在 $S$ 的\Term{open-subset}{开子集} $U$，使得对 $S$ 上任意 $S'$，$H_{S'}$ 形如 ${}_nT_{S'}$，其中 $T_{S'}$ 是 $G_{S'}$ 的某个\Term{maximal-torus}{极大环面}，当且仅当 $S'\to S$ 通过 $U$ \Term{resolution}{分解}。显然可假定 $G$ 的\Term{fiber}{纤维}的\Term{nilpotent-rank}{幂零秩}为常数 $r'$（由 b)，它\Term{locally-constant}{局部常值}）。令 $C=\operatorname{Centr}_G(H)$，这是 $G$ 的一个在 $S$ 上\Term{smooth}{光滑}的\Term{closed-subgroup-prescheme}{闭子群预概形}（XI 6.2）。将 $S$ 换为 $C$ 的\Term{relative-dimension}{相对维数}等于 $r'$ 的点组成的\Term{open-and-closed}{开闭}\Term{subset}{子集}，便可假定 $C$ 处处\Term{relative-dimension}{相对维数}为 $r'$。这时立即可见，$H$ 形如 ${}_nT$（$T$ 是 $G$ 的\Term{maximal-torus}{极大环面}），当且仅当 $C^0$ 有一个在\Term{center}{中心}中的\Term{maximal-torus}{极大环面} $T$，处处\Term{relative-dimension}{相对维数}为 $r$，且 $H={}_nT$。在此判据中将 $S$ 换成 $S$ 上任意 $S'$，就得到待表示的 $S$ 的\Term{subfunctor}{子函子} $U$ 的另一种描述。又由\Term{flat}{平坦}\Term{descent}{下降}，可假定 $G$ 有\Term{maximal-torus}{极大环面} $T_1$。令 $R$ 为 $\operatorname{Transp}_G(T_1,C)$ 的\Term{subfunctor}{子函子} $\operatorname{Transp}_G(T_1,\operatorname{Centr}_{C^0}(C^0))$；后者由 XI 6.2 \Term{representable-functor}{可表示}为 $S$ 上\Term{smooth}{光滑}的\Term{prescheme}{预概形}，而前者由 IX 5.6 a) 立即可知由其中的开\Term{subprescheme}{子预概形}表示，故也在 $S$ 上\Term{smooth}{光滑}。其到 $S$ 的\Term{structure-morphism}{结构态射}因此是开映射，像是\Term{open-subset}{开集}；把 $S$ 换为该像并赋予\Term{induced-structure}{诱导结构}，可假定\Term{structure-morphism}{结构态射}\Term{surjectivity}{满射}。将 1.13 应用于 $C^0$，因 $C^0$ 对于 fpqc \Term{local-scheme}{局部}地有一个在\Term{center}{中心}中的\Term{maximal-torus}{极大环面}，便在整体上也有这样的 $T$，其\Term{relative-dimension}{相对维数}显然等于 $T_1$ 的\Term{relative-dimension}{相对维数}，即 $r$。故对 $H$ 要表达的条件仅为 $H={}_nT$；由 IX 2.10，这仍相当于取 $S$ 的适当开（且闭）\Term{subset}{子集}。
\end{proof}
\begin{proof}[7.1 的证明（续）]
由\Term{lemma}{引理} 7.5，$\mathcal T$ 由 $S$ 上\Term{separated}{分离}\Term{locally-finite-presentation}{局部有限表现}的\Term{prescheme}{预概形}表示，甚至在 $S$ 上\Term{finite-presentation}{有限表现}：只要重看 7.5 的证明，便可确认 $\varphi$ 实际为\Term{quasi-compact}{拟紧}\Term{open-immersion}{开浸入}；或由\Term{faithfully-flat}{忠实平坦}\Term{quasi-compact}{拟紧}\Term{descent}{下降}，归结到 $G$ 有\Term{maximal-torus}{极大环面}的情形，此时 $\mathcal T$ \Term{isomorphism}{同构}于 $G/\operatorname{Norm}_G(T)$。由后一描述，或直接由 XI 2.1，还可见 $\mathcal T$ 在 $S$ 上\Term{smooth}{光滑}。最后，若 $G$ 在 $S$ 上\Term{quasi-projective}{拟射影}，则 $Q$，从而 $\mathcal T$ 亦然。若 $G$ 在 $S$ 上\Term{affine}{仿射}，则 $\mathcal T$ 在 $S$ 上\Term{affine}{仿射}的断言只是重申 5.4 已证的结果。（注意：我不知道，在不假定 $G/S$ \Term{affine}{仿射}时，能否选取 $n$ 使 7.5 的\Term{open-immersion}{开浸入}同时为\Term{closed-immersion}{闭浸入}。）关于 $S$ 为\Term{artinian}{阿廷的}时 $\mathcal T$ 在 $S$ 上\Term{affine}{仿射}的断言，可归结到 $S$ 为域的谱（EGA I 6.1.7），并可假定该域\Term{algebraically-closed}{代数闭}。由下面将证的 f)，只需对 $G/\operatorname{Centr}(G)$ 证明同一断言；
\SourcePage{174}
由 6.1，后者\Term{affine}{仿射}，因而属于前述情形。d) 证毕。

e) 由 IX 6.8，存在 $G'$ 的\Term{subtorus}{子环面} $T'$，使 $u$ 诱导\Term{faithfully-flat}{忠实平坦}\Term{morphism}{态射} $T\to T'$（这将 $T'$ 刻画为 $G'$ 的\Term{subsheaf}{子层} $u(T)$）。设 $C$ 为 $T$ 的\Term{centralizer}{中心化子}，$C'$ 为 $T'$ 的\Term{centralizer}{中心化子}，证明\Term{morphism}{态射} $C\to C'$ \Term{flat}{平坦}，并在 $G'$ 的\Term{fiber}{纤维}\Term{connected}{连通}时\Term{faithfully-flat}{忠实平坦}。$C,C'$ 均在 $S$ 上\Term{flat}{平坦}\Term{finite-presentation}{有限表现}，故可归结到\Term{base}{基底}为域的情形（SGA1 I 5.9），显然还可假定该域\Term{algebraically-closed}{代数闭}。又可将 $G,G'$ 换为 $G^0,G'{}^0$（这不改变 $C^0,C'{}^0$），从而假定 $G,G'$ \Term{connected}{连通}；由 a)，应用 6.6 d) 即可。

f) 由 a)、e)，只需证明关于\Term{cartan-subgroup}{Cartan 子群}的断言。$G$ 的\Term{cartan-subgroup}{Cartan 子群}是某个\Term{maximal-torus}{极大环面}的\Term{centralizer}{中心化子}，故包含 $G$ 的\Term{center}{中心}，更包含 $\operatorname{Ker}u$。因此它形如 $u^{-1}(C')$，其中 $C'=u(C)$ 是 e) 中的 $G'$ 的\Term{cartan-subgroup}{Cartan 子群}。故映射 $C\mapsto u(C)$ \Term{injectivity}{单射}；为证其为\Term{bijection}{双射}，只需证明对 $G'$ 的任意\Term{cartan-subgroup}{Cartan 子群} $C'$，$u^{-1}(C')$ 为 $G$ 的\Term{cartan-subgroup}{Cartan 子群}。问题对 fpqc \Term{local-scheme}{局部}，可假定 $G$ 有\Term{cartan-subgroup}{Cartan 子群} $C_1$。则 $u(C_1)=C'_1$ 是 $G'$ 的\Term{cartan-subgroup}{Cartan 子群}，由 b) 它与 $C'$ 对于 fpqc \Term{local-scheme}{局部}\Term{conjugacy}{共轭}；而 $u^{-1}(C'_1)=C_1$ 是 $G$ 的\Term{cartan-subgroup}{Cartan 子群}，故 $u^{-1}(C')$ 亦然。

也可如下证明 $u^{-1}(C')$ 是\Term{cartan-subgroup}{Cartan 子群}：因 $u$ \Term{flat}{平坦}（SGA1 I 5.9），它在 $S$ 上\Term{flat}{平坦}，于是依定义可归结到\Term{base}{基底}为域的情形，再应用 6.6 e)。
\end{proof}
\begin{SGAresult}{推论}{7.6}
在 7.1 e) 的条件下，$G'$ 对于\Term{etale-topology}{平展拓扑}\Term{local-scheme}{局部}有\Term{maximal-torus}{极大环面}，因而满足对 $G$ 的预备条件。若 $\operatorname{Ker}u$ 还在\Term{center}{中心}中，则 $G$ 的\Term{maximal-torus}{极大环面}\Term{functor}{函子} $\mathcal T_G$ 和\Term{cartan-subgroup}{Cartan 子群}\Term{functor}{函子} $\mathcal C_G$，分别\Term{isomorphism}{同构}于 $G'$ 的相应\Term{functor}{函子} $\mathcal T_{G'}$、$\mathcal C_{G'}$（因而在它们\Term{representable-functor}{可表示}时，由\Term{isomorphism}{同构}的 $S$-\Term{prescheme}{预概形}表示）。
\end{SGAresult}
\textbf{注记 7.7。} a) 与 $G$ 在 $S$ 上\Term{affine}{仿射}时的情况不同（实际上只需 $G$ 的\Term{fiber}{纤维}\Term{affine}{仿射}，见第 XVI 讲），$G$ 的\Term{reductive-rank}{约化秩}\Term{locally-constant}{局部常值}并不\Term{pushout}{推出} $G$ 对于 fpqc \Term{local-scheme}{局部}地有\Term{maximal-torus}{极大环面}。例如，令 $S$ 为\Term{discrete-valuation-ring}{离散赋值环}的谱，$G_1,G_2$ 为 $S$ 上\Term{smooth}{光滑}、\Term{separated}{分离}、\Term{finite-type}{有限型}的\Term{group-scheme}{群概形}，使 $G_1$ 的一般\Term{fiber}{纤维}为\Term{elliptic-curve}{椭圆曲线}，特殊\Term{fiber}{纤维}为 $\mathbb G_m$；$G_2$ 的一般\Term{fiber}{纤维}为 $\mathbb G_m$，特殊\Term{fiber}{纤维}为 $\mathbb G_a$。取 $G=G_1\times_SG_2$。则 $G$ 的两个\Term{fiber}{纤维}的\Term{reductive-rank}{约化秩}都是 $1$，但立即可见 $G$ 对于 fpqc \Term{local-scheme}{局部}也没有\Term{maximal-torus}{极大环面}。另一方面，以下条件很可能足以保证\Term{etale-topology}{平展拓扑}\Term{local-scheme}{局部}存在\Term{maximal-torus}{极大环面}（对\Term{prescheme}{预概形} $S$ 上的\Term{smooth}{光滑}、\Term{separated}{分离}、\Term{finite-type}{有限型}群而言）：$G$ 的\Term{fiber}{纤维}的\Term{reductive-rank}{约化秩}和\Term{abelian-rank}{阿贝尔秩}都是\Term{locally-constant}{局部常值}函数。
\SourcePage{175}
b) 在 7.1 的证明中（尤其是 a)），我们在未假定 $S$ \Term{local-scheme}{局部}\Term{noetherian}{诺特}时引用了 XI 6.2。但在那些应用 XI 6.2 的情形，立即可归结到 $S$ \Term{noetherian}{诺特}且\Term{affine}{仿射}的情况。

以下是 7.1 b) 的一个变体。
\begin{SGAresult}{命题}{7.8}
设 $G$ 是 $S$ 上\Term{smooth}{光滑}\Term{finite-presentation}{有限表现}且\Term{fiber}{纤维}\Term{connected}{连通}的\Term{group-prescheme}{群预概形}，$H$ 是 $S$ 上的\Term{smooth}{光滑}\Term{group-prescheme}{群预概形}，$i:H\to G$ 是 $S$-群的\Term{monomorphism}{单态射}（使 $H$ 为 $G$ 的 $S$-\Term{subgroup}{子群}）。则对 $G$ 的任意\Term{cartan-subgroup}{Cartan 子群}\footnote{证明表明，只需假定 $C$ 是 $G$ 的\Term{smooth}{光滑}\Term{subgroup}{子群}，其每个几何\Term{fiber}{纤维}都是 $G_{\overline s}$ 的某个\Term{maximal-torus}{极大环面}的\Term{subgroup}{子群}的\Term{connected-centralizer}{连通中心化子}。} $C$，$\operatorname{Transp}_G(C,H)$ 由 $G$ 的一个在 $S$ 上\Term{smooth}{光滑}的\Term{closed-subprescheme}{闭子预概形}表示。
\end{SGAresult}
\begin{proof}
由于 $G$ 的\Term{fiber}{纤维}\Term{connected}{连通}，$C$ 亦然（6.6 b)），故\Term{transporter}{输送子}\Term{representable-functor}{可表示}为 $G$ 的\Term{finite-presentation}{有限表现}\Term{closed-subprescheme}{闭子预概形}这一点，包含在 XI 6.11 中。为证其在 $S$ 上\Term{smooth}{光滑}，用标准方法可归结到 $S$ \Term{affine}{仿射}\Term{noetherian}{诺特}，再到 $S$ 为\Term{local}{局部}\Term{artinian}{阿廷的}，又由\Term{descent}{下降}归结到 $S$ 的\Term{residue-field}{剩余域}\Term{algebraically-closed}{代数闭}。这时 $C$ 有\Term{maximal-torus}{极大环面} $T$，它也是 $G$ 的\Term{maximal-torus}{极大环面}。可假定\Term{fiber}{纤维} $H_0$ 的\Term{reductive-rank}{约化秩}和\Term{nilpotent-rank}{幂零秩}等于 $G_0$ 的相应秩（否则\Term{transporter}{输送子}为空）。于是利用 $C$ 的\Term{connected}{连通性}及 $H$ 的\Term{maximal-torus}{极大环面}在 $H$ 中的\Term{centralizer}{中心化子}\Term{smooth}{光滑}这一事实，立即可见
\[
\operatorname{Transp}_G(C,H)=\operatorname{Transp}_G(T,H).
\]
已知右端\Term{smooth}{光滑}（XI 2.5），故左端亦然。上述论证更一般地证明了下述命题的 b)。
\end{proof}
\begin{SGAresult}{命题}{7.9}
设 $G$ 和 $i:H\to G$ 如 7.8，再假定对每个 $s\in S$，$H_s$ \Term{connected}{连通}，且与 $G_s$ 有相同的\Term{reductive-rank}{约化秩}及\Term{nilpotent-rank}{幂零秩}（即 $H_s$ 包含 $G_s$ 的一个\Term{cartan-subgroup}{Cartan 子群}）。则：

 a) $\operatorname{Norm}_G(H)$ 由 $G$ 的一个在 $S$ 上\Term{finite-presentation}{有限表现}的\Term{closed-subprescheme}{闭子预概形} $\operatorname{Norm}_G(H)$ 表示，\Term{canonical}{典范}\Term{monomorphism}{单态射} $H\to\operatorname{Norm}_G(H)$ 是\Term{open-immersion}{开浸入}；因而 $i$ 是\Term{immersion}{浸入}，并有 $H=\operatorname{Norm}_G(H)^0$。

 b) 对 $G$ 的任意\Term{cartan-subgroup}{Cartan 子群} $C$，$\operatorname{Transp}_G(C,H)$ 是 $G$ 的\Term{closed-subprescheme}{闭子预概形}，在 $S$ 上\Term{smooth}{光滑}，且\Term{structure-morphism}{结构态射}\Term{surjectivity}{满射}。若 $C$ 是 $G$ 的某个\Term{maximal-torus}{极大环面} $T$ 的\Term{centralizer}{中心化子}，还有
\[
\operatorname{Transp}_G(T,H)=\operatorname{Transp}_G(C,H).
\]
 c) 设 $C$ 是 $H$ 的\Term{subgroup-prescheme}{子群预概形}。则 $C$ 是 $H$ 的\Term{cartan-subgroup}{Cartan 子群}，当且仅当它是 $G$ 的\Term{cartan-subgroup}{Cartan 子群}。

 d) 假定 $G$ 对于\Term{etale-topology}{平展拓扑}或 \Term{fpqc-topology}{fpqc 拓扑}\Term{local-scheme}{局部}地有\Term{cartan-subgroup}{Cartan 子群}（相应地，\Term{maximal-torus}{极大环面}），则 $H$ 亦然。
\end{SGAresult}
\SourcePage{176}
\begin{proof}
a) $\operatorname{Norm}_G(H)$ \Term{representable-functor}{可表示}为 $G$ 的一个在 $S$ 上\Term{finite-presentation}{有限表现}的\Term{closed-subprescheme}{闭子预概形}，这包含在 XI 6.11 中。由于 $H$ 在 $S$ 上\Term{smooth}{光滑}，因而\Term{flat}{平坦}且\Term{locally-finite-presentation}{局部有限表现}，为验证 $H\to\operatorname{Norm}_G(H)$ 是\Term{open-immersion}{开浸入}，只需在\Term{fiber}{纤维}上验证（VI$_{\mathrm B}$ 2.6），从而归结到 $S$ 为\Term{algebraically-closed-field}{代数闭域} $k$ 的谱。此时（第 VI 讲\footnote{编者注：未在第 VI$_{\mathrm B}$ 讲找到此引用，参见 M. Demazure 与 P. Gabriel，\textit{Groupes algébriques}，I，Masson（1970），推论 II.5.6。}）只需验证相应\Term{lie-algebra}{李代数}上的\Term{homomorphism}{同态}为\Term{isomorphism}{同构}，亦即 $(\mathfrak g/\mathfrak h)^H=0$，其中 $\mathfrak g,\mathfrak h$ 为 $G,H$ 的\Term{lie-algebra}{李代数}，上标 $H$ 表示 $H$ 下的\Term{invariant}{不变量}（参见 II 5.2.3 (i)）。根据假设，$H$ 包含 $G$ 的某个\Term{cartan-subgroup}{Cartan 子群} $C$，它是 $G$ 的\Term{maximal-torus}{极大环面} $T$ 的\Term{centralizer}{中心化子}，故只需证明
\[
(\mathfrak g/\mathfrak h)^T=0.
\]
利用 $T$ 的表示完全可约（I 4.7.3），这来自相应关系 $(\mathfrak g/\mathfrak c)^T$\SGAsourcegap{原书此处仅印出\Term{invariants}{不变量}式，未印出该关系的其余部分。}，其中 $\mathfrak c=\operatorname{Lie}(C)$。后一个关系\Term{equivalence}{等价}于
\[
\mathfrak g^T=\mathfrak c,
\]
它意味着 $T$ 的\Term{centralizer}{中心化子} $C$ 和\Term{normalizer}{正规化子} $N$ 有相同的\Term{lie-algebra}{李代数}；这是因为 $C$ 是 $N$ 的\Term{open-subgroup}{开子群}（XI 5.9）。a) 证毕。

b) 如上所述，证明已在 7.8 给出。

c) 因 $H$ 是 $G$ 的\Term{subprescheme}{子预概形}，该断言显然可归结到 $S$ 为\Term{algebraically-closed-field}{代数闭域}的谱，此时由对 $H$ 的假设立即得出。

d) 使用 b)、c) 及“Hensel \Term{lemma}{引理}”XI 1.10。
\end{proof}
\begin{SGAresult}{推论}{7.10}
设 $G,H$ 如 7.9，并假定对 $S$ 上每个\Term{algebraically-closed-field}{代数闭域} $k$，$G_k(k)$ 中\Term{normalization}{正规化} $H_k$ 的每个元素都属于 $H_k(k)$。则 $H$ 是 $G$ 的\Term{closed-subprescheme}{闭子预概形}，并等于其自身的\Term{normalizer}{正规化子}。
\end{SGAresult}
这由 7.9 a) 显然得出。我们将把 7.10 特别应用于 $G$ 的\Term{borel-subgroup}{Borel 子群}（更一般地，\Term{parabolic-subgroup}{抛物子群}）。
\begin{SGAresult}{推论}{7.11}
设 $G,H$ 如 7.9。则 $H$ 包含 $G$ 的每个\Term{subgroup-prescheme}{子群预概形} $Z$，只要 $Z$ 在 $G$ 的\Term{center}{中心}中且在 $S$ 上\Term{flat}{平坦}\Term{finite-presentation}{有限表现}。
\end{SGAresult}
\begin{proof}
如通常一样，可归结到 $S$ \Term{affine}{仿射}\Term{noetherian}{诺特}，再到 $S$ 为\Term{artinian}{阿廷的}，从而处于 7.1 的条件下。由 7.9 c)，可假定 $H$ 包含 $G$ 的\Term{cartan-subgroup}{Cartan 子群} $C$，只需证明 $C\supset Z$。因 $S$ 为\Term{artinian}{阿廷的}，$G'=G/Z$ 由 $S$ 上\Term{finite-type}{有限型}\Term{group-prescheme}{群预概形}表示，\Term{canonical}{典范}\Term{morphism}{态射} $u:G\to G'$ \Term{faithfully-flat}{忠实平坦}且\Term{kernel}{核}为 $Z$（VI$_{\mathrm A}$ 3.2）。$G'$ 显然在 $S$ 上\Term{smooth}{光滑}，可应用 7.1 f)，得到 $C$ 形如 $u^{-1}(C')$，故包含 $Z$。
\end{proof}
\SourcePage{177}
\begin{SGAresult}{推论}{7.12}
设 $G,G'$ 是 $S$ 上\Term{smooth}{光滑}\Term{finite-presentation}{有限表现}的两个\Term{group-prescheme}{群预概形}，$u:G\to G'$ 是\Term{faithfully-flat}{忠实平坦}群\Term{homomorphism}{同态}（即对每个 $s\in S$，$u_s$ \Term{faithfully-flat}{忠实平坦}）。假定 $\operatorname{Ker}u$ 在\Term{center}{中心}中且 $G$ 的\Term{fiber}{纤维}\Term{connected}{连通}。则映射
\[
H'\longmapsto H=u^{-1}(H')
\]
给出如下两\Term{set}{集合}间的一一对应：一边为 $G'$ 的\Term{subgroup-prescheme}{子群预概形} $H'$，在 $S$ 上\Term{smooth}{光滑}\Term{finite-presentation}{有限表现}，且在每个 $s\in S$ 处与 $G'$ 有相同的\Term{reductive-rank}{约化秩}和\Term{nilpotent-rank}{幂零秩}；另一边为 $G$ 的\Term{subgroup-prescheme}{子群预概形} $H$，在 $S$ 上\Term{smooth}{光滑}\Term{finite-presentation}{有限表现}，且在每个 $s\in S$ 处与 $G$ 有相同的\Term{reductive-rank}{约化秩}和\Term{nilpotent-rank}{幂零秩}。$H$ 的\Term{fiber}{纤维}\Term{connected}{连通}，当且仅当 $H'$ 的\Term{fiber}{纤维}\Term{connected}{连通}。
\end{SGAresult}
\begin{proof}
设 $H$ 是 $G$ 的具有上述性质的\Term{subgroup-prescheme}{子群预概形}。由 7.10，$H$ 包含 $Z=\operatorname{Ker}u$，因此由\Term{faithfully-flat}{忠实平坦}\Term{descent}{下降}理论，$H=u^{-1}(H')$，其中 $H'$ 是 $G'$ 的某个唯一确定的\Term{subgroup-prescheme}{子群预概形}；由 6.6 e) 立即可见，后者具有上述性质。又若 $H$ 的\Term{fiber}{纤维}\Term{connected}{连通}，则 $H'=H/Z$ 显然亦然。因此只需证明：若从具有所述性质的 $G'$ 的\Term{subgroup}{子群} $H'$ 出发，则 $u^{-1}(H')=H$ 在 $G$ 中具有同样性质；且若 $H'$ 的\Term{fiber}{纤维}\Term{connected}{连通}，则 $H$ 亦然。由 6.6 e)，只需证明 $H$ 在 $S$ 上\Term{smooth}{光滑}（相应地，且\Term{fiber}{纤维}\Term{connected}{连通}）。$H$ 是在 $S$ 上\Term{flat}{平坦}的 $H'$ 沿\Term{flat}{平坦}\Term{morphism}{态射} $u$ 的\Term{inverse-image}{逆像}，故已在 $S$ 上\Term{flat}{平坦}；只需验证 $H$ 的几何\Term{fiber}{纤维}\Term{smooth}{光滑}（相应地，且\Term{connected}{连通}），这便归结到 $S$ 为\Term{algebraically-closed-field}{代数闭域}的谱。此时 $H'$ 包含 $G'$ 的\Term{cartan-subgroup}{Cartan 子群} $C'$，故 $H$ 包含 $C'$ 的\Term{inverse-image}{逆像} $C$；由 6.6 e)，$C$ 是 $G$ 的\Term{cartan-subgroup}{Cartan 子群}，从而\Term{smooth}{光滑}\Term{connected}{连通}。因此 7.11 由下述\Term{lemma}{引理}得出。
\end{proof}
\begin{SGAresult}{\Term{lemma}{引理}}{7.13}
设 $G,G'$ 是 $S$ 上\Term{flat}{平坦}\Term{finite-presentation}{有限表现}的两个\Term{group-prescheme}{群预概形}，$u:G\to G'$ 是\Term{faithfully-flat}{忠实平坦}群\Term{homomorphism}{同态}，$C'$ 是 $G$ 的\Term{subgroup-prescheme}{子群预概形}，在 $S$ 上\Term{finite-presentation}{有限表现}，使 $C=u^{-1}(C')$ 在 $S$ 上\Term{smooth}{光滑}（相应地，\Term{fiber}{纤维}\Term{connected}{连通}）。则对 $G'$ 的任意在 $S$ 上\Term{finite-presentation}{有限表现}、包含 $C'$ 且在 $S$ 上\Term{smooth}{光滑}（相应地，\Term{fiber}{纤维}\Term{connected}{连通}）的\Term{subgroup-prescheme}{子群预概形} $H'$，其\Term{inverse-image}{逆像} $H=u^{-1}(H')$ 在 $S$ 上\Term{smooth}{光滑}（相应地，\Term{fiber}{纤维}\Term{connected}{连通}）。
\end{SGAresult}
\begin{proof}
如上所述，该命题立即可归结到 $S$ 为域的谱。此时 $H$ 是以 $H/C$ 为底、以 $C$ 为群的\Term{principal-bundle}{主丛}（第 VI$_{\mathrm B}$ 讲第 9 节）；又 $H/C\simeq H'/C'$（参见第 IV 讲）。$H'$ \Term{smooth}{光滑}（相应地，\Term{connected}{连通}），故 $H'/C'$，从而 $H/C$ 亦然。依假设 $C$ 也如此，于是\Term{bundle}{丛} $H$ 立即也如此。
\end{proof}

\SGAsection{8}{不必仿射的群概形中的半单元素以及极大环面的并与交}
本节始终以 $G$ 表示 $S$ 上\Term{smooth}{光滑}\Term{finite-presentation}{有限表现}且\Term{fiber}{纤维}\Term{connected}{连通}的\Term{group-prescheme}{群预概形}。
\SourcePage{178}
先设 $S$ 为\Term{algebraically-closed-field}{代数闭域}的谱。当 $G$ \Term{affine}{仿射}时，BIBLE 4 第 4 节中已经定义了 $G(k)$ 的\Term{semisimple-element}{半单元素}概念；立即可见，该概念在\Term{base-field}{基域}的\Term{algebraically-closed}{代数闭}\Term{group-extension}{扩张} $k'/k$ 下\Term{invariant}{不变}。又由 BIBLE 6 \Term{theorem}{定理} 5 (c)，$g\in G(k)$ \Term{semisimple-element}{半单}，当且仅当它包含在 $G$ 的某个\Term{maximal-torus}{极大环面}内。不再假定 $G$ \Term{affine}{仿射}时，由 Chevalley \Term{theorem}{定理}，$G$ \Term{canonical}{典范}地写成\Term{abelian-variety}{阿贝尔簇}被\Term{smooth}{光滑}\Term{connected}{连通}\Term{affine}{仿射}\Term{algebraic-group}{代数群}的\Term{extension}{扩张}：
\[
e\longrightarrow V\longrightarrow G\longrightarrow A\longrightarrow e.\tag{*}
\]
称 $G(k)$ 的元素 $g$ 为\Term{semisimple-element}{半单的}，若它是 $V(k)$ 的\Term{semisimple-element}{半单元素}。$V$ 的\Term{maximal-torus}{极大环面}显然就是 $G$ 的\Term{maximal-torus}{极大环面}，故也\Term{equivalence}{等价}于说 $g$ 属于 $G$ 的某个\Term{maximal-torus}{极大环面}。这一概念显然仍在\Term{base-field}{基域} $k$ 的\Term{algebraically-closed}{代数闭}\Term{group-extension}{扩张} $k'/k$ 下\Term{invariant}{不变}。

现在设 $S$ 为任意域 $k$ 的谱，$g\in G$。选择 $\kappa(g)$ 的\Term{algebraically-closed}{代数闭}\Term{group-extension}{扩张} $K$，可见 $g$ 是 $G$ 的一个\Term{geometric-point}{几何点} $g'$ 的像，$g'$ 取值于 $k$ 的\Term{algebraically-closed}{代数闭}\Term{group-extension}{扩张} $K$ 中。若 $g'$ \Term{semisimple-element}{半单}，便称 $g$ \Term{semisimple-element}{半单}；由上面的论述，这与 $g'$ 的具体选取无关。若 $k'$ 为 $k$ 的\Term{group-extension}{扩张}，则 $G_{k'}$ 的\Term{semisimple-element}{半单元素}\Term{set}{集合}为 $G$ 的\Term{semisimple-element}{半单元素}\Term{set}{集合} $G^{\mathrm{ss}}$ 的\Term{inverse-image}{逆像}。

最后设 $S$ 任意。点 $g\in G$ 称为\Term{semisimple-element}{半单的}，若它在所属的\Term{fiber}{纤维} $G_s$ 中\Term{semisimple-element}{半单}。故对任意\Term{morphism}{态射} $S'\to S$，$G_{S'}^{\mathrm{ss}}$ 是 $G^{\mathrm{ss}}$ 的\Term{inverse-image}{逆像}。假定 1.10 中定义的\Term{functor}{函子} $\mathcal T$（\Term{maximal-torus}{极大环面函子}）由 $S$ 上的\Term{finite-presentation}{有限表现}\Term{prescheme}{预概形}表示（例如，若 $G$ 对于 fpqc \Term{local-scheme}{局部}地有\Term{maximal-torus}{极大环面}，由 7.1 d)，至少当 $G$ 在 $S$ 上\Term{quasi-projective}{拟射影}时如此）。考虑 $G_{\mathcal T}$ 的\Term{canonical}{典范}\Term{maximal-torus}{极大环面} $\mathbb T$（“$G$ 的泛\Term{maximal-torus}{极大环面}”），及\Term{projection}{投影} $G_{\mathcal T}\to G$ 诱导的\Term{morphism}{态射}
\[
u:\mathbb T\longrightarrow G.
\]
由定义立即可见，$G^{\mathrm{ss}}$ 正是上述\Term{morphism}{态射}的像。我们将由此\Term{pushout}{推出}：
\begin{SGAresult}{命题}{8.1}
$G$ 的\Term{semisimple-element}{半单元素}\Term{set}{集合} $G^{\mathrm{ss}}$ 是\Term{locally-constructible}{局部可构造的}（因此若 $S$，从而 $G$ \Term{quasi-compact}{拟紧}且\Term{quasi-separated}{拟分离}，则它\Term{constructible}{可构造}）。
\end{SGAresult}
\begin{proof}
如通常一样，归结到 $S$ \Term{affine}{仿射}\Term{noetherian}{诺特}；再由通常的\Term{noetherian}{诺特}\Term{constructibility}{可构造性}判据（EGA $0_{\mathrm{III}}$ 9.2.3），可假定 $S$ \Term{integral}{整}，只需证明存在 $S$ 的\Term{nonempty}{非空}\Term{open-subset}{开集} $U$，使 $G^{\mathrm{ss}}|U$ \Term{constructible}{可构造}。将 $U$ 取得充分小，并在必要时换为一个\Term{finite-cover}{有限覆盖}，可假定 $G$ 在 $S$ 上\Term{separated}{分离}且包含\Term{maximal-torus}{极大环面} $T$。由 7.1 d)，\Term{functor}{函子} $\mathcal T$ 由 $S$ 上\Term{finite-presentation}{有限表现}的\Term{prescheme}{预概形}表示，$\mathbb T$ 亦然，所以它在 $G$ 中的像\Term{constructible}{可构造}。
\end{proof}
\SourcePage{179}
重新假定 $k$ 为域，并考虑由上述\Term{morphism}{态射}生成的 $G$ 的\Term{subgroup}{子群} $H$（参见 VI$_{\mathrm B}$ 1.2）。它是 $G$ 的\Term{smooth}{光滑}\Term{subgroup-scheme}{子群概形}，因 $\mathbb T$ \Term{connected}{连通}而\Term{connected}{连通}，且其构造显然与\Term{base-field}{基域}的任意\Term{group-extension}{扩张}相容（参见第 VI 讲），因为 $\mathbb T$ 的构造如此。$k$ \Term{algebraically-closed}{代数闭}时，立即可见 $H$ 也是由 $G$ 的\Term{maximal-torus}{极大环面}，或\Term{equivalence}{等价}地由 $G$ 的\Term{torus}{环面}生成的\Term{algebraic-subgroup}{代数子群}；它还是包含 $G(k)$ 的全部\Term{semisimple-element}{半单元素}的最小\Term{algebraic-subgroup}{代数子群}。（实际上，只要 $k$ 为无限域，这些对 $H$ 的刻画仍成立，因为第 XIV 讲将证明，$\mathcal T$ 的 $k$-\Term{rational-point}{有理点}集在 $\mathcal T$ 中\Term{dense}{稠密}。）此外，$H$ 是 $G$ 的\Term{normal-subgroup}{正规子群}：为此可假定 $k$ \Term{algebraically-closed}{代数闭}；因 $H,G$ 在 $k$ 上\Term{smooth}{光滑}，只需检验 $H$ 对 $\operatorname{int}(g)$（$g\in G(k)$）这些\Term{inner-automorphism}{内自同构}\Term{stable}{稳定}，这是显然的。（甚至可证明 $H$ 是 $G$ 的\Term{characteristic-subgroup}{特征子群}，即在 $\operatorname{Aut}_{k\text{-gr}}(G)$ 下\Term{stable}{稳定}。）由 6.6 d) 立即可知，$G/H$ 的\Term{reductive-rank}{约化秩}为零。更准确地说，若 $K$ 是 $G$ 的\Term{normal-subgroup}{正规代数子群}，则由 $k$ \Term{algebraically-closed}{代数闭}时 $G/K$ 的\Term{maximal-torus}{极大环面}正是 $G$ 的\Term{maximal-torus}{极大环面}的\Term{direct-image}{正像}，立即可知 $G/K$ 的\Term{reductive-rank}{约化秩}为零，当且仅当 $K$ 包含 $G$ 的所有\Term{maximal-torus}{极大环面}（假定 $k$ \Term{algebraically-closed}{代数闭}），亦即当且仅当 $K$ 包含 $H$（$k$ 任意）。因此，$H$ 是使 $G/H$ 的\Term{reductive-rank}{约化秩}为零的最小\Term{normal-subgroup}{正规代数子群}。$H$ 还有另一个显然的刻画：它是与 $G$ 有相同\Term{reductive-rank}{约化秩}的最小\Term{algebraic-subgroup}{代数子群}。最后，$H$ 是\Term{affine}{仿射的}：仍可假定 $k$ \Term{algebraically-closed}{代数闭}，回到本节开头的\Term{exact-sequence}{正合列} $(*)$，$G$ 的\Term{maximal-torus}{极大环面}包含在 $V$ 中，故 $H$ 亦然；$V$ \Term{affine}{仿射}，故 $H$ \Term{affine}{仿射}。归纳上述结果如下。
\begin{SGAresult}{命题}{8.2}
设 $G$ 是域 $k$ 上\Term{smooth}{光滑}\Term{connected}{连通}的\Term{algebraic-group}{代数群}，$\overline k$ 为 $k$ 的\Term{algebraic-closure}{代数闭包}。存在 $G$ 的\Term{algebraic-subgroup}{代数子群} $H$，使 $H_{\overline k}$ 为 $G_{\overline k}$ 的\Term{maximal-torus}{极大环面}（或\Term{torus}{环面}）生成的\Term{algebraic-subgroup}{代数子群}。$H$ 也可刻画为与 $G$ 有相同\Term{reductive-rank}{约化秩}的最小\Term{algebraic-subgroup}{代数子群}，或使 $G/H$ 的\Term{reductive-rank}{约化秩}为零的最小\Term{normal-subgroup}{正规代数子群}。它是 $G$ 的\Term{smooth}{光滑}\Term{connected}{连通}\Term{normal-subgroup}{正规}\Term{affine}{仿射}\Term{subgroup}{子群}，其构造与\Term{base-field}{基域}的任意\Term{group-extension}{扩张}交换。
\end{SGAresult}
为使用借助 $G/H$ 给出的 $H$ 的刻画，宜将下述推论明确写出。
\begin{SGAresult}{推论}{8.3}
设 $G$ 是\Term{algebraically-closed-field}{代数闭域} $k$ 上\Term{smooth}{光滑}\Term{connected}{连通}的\Term{algebraic-group}{代数群}。$G$ 的\Term{reductive-rank}{约化秩}为零（即用 8.2 的记号，$H=(e)$），当且仅当 $G$ 是\Term{abelian-variety}{阿贝尔簇}被\Term{smooth}{光滑}\Term{connected}{连通}\Term{unipotent}{幺幂}\Term{algebraic-group}{代数群}的\Term{extension}{扩张}（即由\Term{isomorphism}{同构}于\Term{additive-group}{加法群} $\mathbb G_a$ 的群逐次\Term{extension}{扩张}而成）。
\end{SGAresult}
\begin{proof}
由 Chevalley \Term{exact-sequence}{正合列} $(*)$，只需证明\Term{smooth}{光滑}\Term{connected}{连通}\Term{affine}{仿射}群 $V$ 的\Term{reductive-rank}{约化秩}为零，当且仅当 $V$ \Term{unipotent}{幺幂}；
\SourcePage{180}
这包含在 BIBLE 6.4 \Term{theorem}{定理} 4 推论 3 中。
\end{proof}
因此，用 8.2 的记号，$H$ 是使 $(G/H)_{\overline k}$ 为\Term{abelian-variety}{阿贝尔簇}被\Term{smooth}{光滑}\Term{connected}{连通}\Term{affine}{仿射}\Term{unipotent}{幺幂}\Term{algebraic-group}{代数群}的\Term{extension}{扩张}的最小\Term{normal-subgroup}{正规代数子群}。还可得出：
\begin{SGAresult}{推论}{8.4}
用 8.2 的记号，$G=H$（即 $G_{\overline k}$ 由其\Term{maximal-torus}{极大环面}生成），当且仅当 $G$ \Term{affine}{仿射}，且从 $G_{\overline k}$ 到\Term{additive-group}{加法群}的任意\Term{homomorphism}{同态}都平凡。
\end{SGAresult}
\textbf{注记 8.5。} a) 设 $\overline V$ 为 $G_{\overline k}$ 的最大\Term{smooth}{光滑}\Term{connected}{连通}\Term{affine}{仿射}\Term{algebraic-subgroup}{代数子群}（因而 $G_{\overline k}$ 为\Term{abelian-variety}{阿贝尔簇}被 $\overline V$ 的\Term{extension}{扩张}）。众所周知（Rosenlicht\footnote{编者注：此处应给出参考文献。\SGAsourcegap{原书此处未给出 Rosenlicht 结果的参考文献。}}），若 $k$ 不\Term{perfect}{完美}，$\overline V$ 一般并不“定义在 $k$ 上”，即一般不存在 $G$ 的\Term{algebraic-subgroup}{代数子群} $V$ 使 $V_{\overline k}=\overline V$。但若 $\overline V$ 由其\Term{maximal-torus}{极大环面}生成，即 $\overline V$ 不存在\Term{isomorphism}{同构}于\Term{additive-group}{加法群} $\mathbb G_{a,\overline k}$ 的\Term{quotient-group}{商群}，则这样的 $V$ 存在，就是 8.2 中的群 $H$。由此可见，在这一有理性问题中，与许多其他问题一样（例如参见第 XIV 讲第 6 节），所有困难都来自\Term{unipotent}{幺幂}群，即\Term{additive-group}{加法群}；而（足够多的）\Term{multiplicative-type}{乘法型}群的存在则保证事情顺利。

b) 用 8.2 的记号，还可引入 $G/H$ 的\Term{commutator-subgroup}{换位子群}在 $G$ 中的\Term{inverse-image}{逆像} $H'$。则 $H'$ 是使 $(G/H')_{\overline k}$ 为\Term{abelian-variety}{阿贝尔簇}被\Term{smooth}{光滑}\Term{connected}{连通}\Term{affine}{仿射}\Term{unipotent}{幺幂}\Term{algebraic-group}{代数群}的交换\Term{extension}{扩张}的最小\Term{normal-subgroup}{正规代数子群}。$H'$ 仍是 $G$ 的\Term{smooth}{光滑}\Term{connected}{连通}\Term{affine}{仿射}\Term{subgroup}{子群}。设 $H''$ 为 $p^n(G/H')$ 在 $G$ 中的\Term{inverse-image}{逆像}，其中 $n$ 充分大，$p$ 为\Term{characteristic}{特征}（假定 $p>0$）。容易看出 $H''/H'$ 为\Term{abelian-variety}{阿贝尔簇}，而 $H''$ 是使 $G/H''$ 为交换\Term{smooth}{光滑}\Term{connected}{连通}\Term{affine}{仿射}\Term{unipotent}{幺幂}\Term{algebraic-group}{代数群}的最小\Term{normal-subgroup}{正规代数子群}。于是容易看出，$\overline V=H_{\overline k}$（a) 的记号），即从 $\overline V$ 到\Term{additive-group}{加法群}的任意\Term{homomorphism}{同态}都平凡，当且仅当 $H''=G$，亦即当且仅当从 $G_{\overline k}$ 到\Term{additive-group}{加法群}的任意\Term{homomorphism}{同态}都平凡。

最后，仿照 4.10，将第 4 节中建立的\Term{reductive-center}{约化中心}概念推广到\Term{fiber}{纤维}\Term{connected}{连通}的\Term{smooth}{光滑}群。令 $Z$ 为 $G$ 的以下\Term{subfunctor}{子函子}：
\[
\begin{split}
Z(S')=\{\,&G_{S'}\text{ 在 }S'\text{ 上的截面 }g'\mid
\text{对 }S'\text{ 上任意 }S''\text{ 及 }G_{S''}\text{ 的任意极大环面 }T'',\\
&g'\text{ 沿 }S''\to S'\text{ 的逆像 }g''\text{ 是 }T''\text{ 在 }S''\text{ 上的截面}\,\}.
\end{split}
\]
引入\Term{functor}{函子} $\mathcal T(S')=\{G_{S'}\text{ 的极大环面}\}$ 及 $\mathbb T(S')=\{(T',g')\}$，其中 $T'$ 是 $G_{S'}$ 的\Term{maximal-torus}{极大环面}，$g'$ 是 $T'$ 在 $S'$ 上的\Term{section}{截面}。可见 $\mathbb T$ 为 $\mathcal T\times_SG=\mathcal T_G$ 的\Term{subfunctor}{子函子}，并可写成
\[
Z=\prod_{\mathcal T_G/G}\mathbb T/\mathcal T_G.
\]
\SourcePage{181}
利用 XI 6.8，以及在某些条件下保证 $\mathcal T$ 可由 $S$ 上\Term{smooth}{光滑}\Term{prescheme}{预概形}表示的 7.1 d)，可在某些条件下\Term{pushout}{推出} $Z$ \Term{representable-functor}{可表示}；不过下面将用更直接的方法取得这些结论。
\begin{SGAresult}{定义}{8.6}
设 $G$ 为 $S$ 上\Term{smooth}{光滑}\Term{finite-presentation}{有限表现}且\Term{fiber}{纤维}\Term{connected}{连通}的\Term{group-prescheme}{群预概形}。称 $G$ 有\Term{reductive-center}{约化中心}，若上述\Term{functor}{函子} $Z$（显然为 $G$ 的\Term{subgroup}{子群}）由\Term{multiplicative-type}{乘法型}群表示。这时称 $Z$ 为 $G$ 的\Term{reductive-center}{约化中心}。
\end{SGAresult}
注意，若 $Z$ 是 $G$ 的\Term{reductive-center}{约化中心}，则对任意\Term{base-change}{基变换} $S'\to S$，$Z_{S'}$ 是 $G_{S'}$ 的\Term{reductive-center}{约化中心}；另一方面，\Term{reductive-center}{约化中心}的存在性显然对于 \Term{fpqc-topology}{fpqc 拓扑}是\Term{local-scheme}{局部}问题。至于“\Term{reductive-center}{约化中心}”这一名称，注意 $Z$ 总在\Term{center}{中心}中，因为显然 $Z$ 在 $\operatorname{Aut}_S(G)$ 下\Term{invariant}{不变}，因而是 $G$ 的\Term{normal-subgroup}{正规子群}，再应用 IX 5.5 即可。

这里，\Term{lemma}{引理} 4.5 应替换为：
\begin{SGAresult}{\Term{lemma}{引理}}{8.7}
设 $u:H\to G$ 是\Term{central-homomorphism}{中心同态}，$H$ 在 $S$ 上为\Term{multiplicative-type}{乘法型}且\Term{finite-type}{有限型}。假定对 $S$ 上任意\Term{algebraically-closed-field}{代数闭域} $k$，$u_k(H_k)$ 包含在 $G_k$ 的最大\Term{affine}{仿射}\Term{smooth}{光滑}\Term{connected}{连通}\Term{subgroup}{子群}内。则 $u$ 通过 $G$ 的任意\Term{maximal-torus}{极大环面} $T$ \Term{resolution}{分解}（所以实际上 $u$ 通过前面定义的 $G$ 的\Term{subfunctor}{子函子} $Z$ \Term{resolution}{分解}）。
\end{SGAresult}
\begin{proof}
使用 IX 5.1 bis 的一个简单变体（将等号换成包含号），可归结到 $S$ 为域 $k$ 的谱，并可假定 $k$ \Term{algebraically-closed}{代数闭}。（先归结到 $S$ \Term{affine}{仿射}\Term{noetherian}{诺特}，再到\Term{artinian}{阿廷}情形，然后用 IX 3.6。）$T$ 包含在 $G$ 的最大\Term{affine}{仿射}\Term{smooth}{光滑}\Term{connected}{连通}\Term{subgroup}{子群} $V$ 内，因而只需考虑 $G=V$，即 $G$ \Term{affine}{仿射}的情形；此时结论已在 4.5 中证明。
\end{proof}
\begin{SGAresult}{命题}{8.8}
设 $G$ 为 $S$ 上\Term{smooth}{光滑}\Term{finite-presentation}{有限表现}且\Term{fiber}{纤维}\Term{connected}{连通}的\Term{group-prescheme}{群预概形}。

 a) 若 $S$ 是域 $k$ 的谱，则 $G$ 有\Term{reductive-center}{约化中心} $Z$。当 $G$ 为\Term{abelian-variety}{阿贝尔簇}被\Term{smooth}{光滑}\Term{connected}{连通}\Term{affine}{仿射}\Term{algebraic-group}{代数群} $V$ 的\Term{extension}{扩张}时（例如 $k$ \Term{algebraically-closed}{代数闭}时），$Z$ 也是 $V$ 的\Term{reductive-center}{约化中心}，且为 $V$ 的最大\Term{central-subgroup}{中心}\Term{multiplicative-type}{乘法型子群}。

 b) 设 $Z$ 为 $G$ 的\Term{multiplicative-type}{乘法型}\Term{subgroup-prescheme}{子群预概形}。则 $Z$ 为 $G$ 的\Term{reductive-center}{约化中心}，当且仅当对每个 $s\in S$，$Z_s$ 为 $G_s$ 的\Term{reductive-center}{约化中心}。这时，$Z$ 是 $G$ 的满足下述条件的最大\Term{multiplicative-type}{乘法型子群} $K$：对每个 $s\in S$，$K_s$ 包含在 $G_s$ 的\Term{reductive-center}{约化中心}中；更一般地，对任意\Term{homomorphism}{同态} $u:H\to G$，其中 $H$ 在 $S$ 上为\Term{multiplicative-type}{乘法型}且\Term{finite-type}{有限型}，若对 $S$ 上任意\Term{algebraically-closed-field}{代数闭域} $k$，$u_k:H_k\to G_k$ 通过 $G_k$ 的最大\Term{smooth}{光滑}\Term{connected}{连通}\Term{affine}{仿射}\Term{subgroup}{子群}\Term{resolution}{分解}，则 $u$ 通过 $Z$ \Term{resolution}{分解}（特别地，$u$ 为\Term{central-homomorphism}{中心同态}）。
\SourcePage{182}
 c) 若 $G$ 对于 \Term{fpqc-topology}{fpqc 拓扑}\Term{local-scheme}{局部}地有\Term{maximal-torus}{极大环面}，则 $G$ 有\Term{reductive-center}{约化中心}。

 d) 设 $T$ 为 $G$ 的\Term{maximal-torus}{极大环面}。则 $T\cap\operatorname{Centr}(G)=\operatorname{Ker}(T\to\operatorname{GL}(\mathfrak g))$，且它是 $G$ 的\Term{reductive-center}{约化中心}。

 e) 设 $Z$ 为 $G$ 的\Term{reductive-center}{约化中心}，并假定 $G'=G/Z$ \Term{representable-functor}{可表示}（例如 $S$ 为\Term{artinian}{阿廷的}）。则 $G'$ 以单位\Term{subgroup}{子群}为\Term{reductive-center}{约化中心}，而 $T'\mapsto T=u^{-1}(T')$ 给出 $G'$ 的\Term{maximal-torus}{极大环面}与 $G$ 的\Term{maximal-torus}{极大环面}之间的一一对应。
\end{SGAresult}
\begin{proof}
a) 设 $S$ 为域 $k$ 的谱。为证\Term{reductive-center}{约化中心}存在，可假定 $k$ \Term{algebraically-closed}{代数闭}，从而归结到 $G$ 为\Term{abelian-variety}{阿贝尔簇}被\Term{smooth}{光滑}\Term{connected}{连通}\Term{affine}{仿射}\Term{algebraic-group}{代数群} $V$ 的\Term{extension}{扩张}的情形。对 $k$ 上任意 $S'$，$G_{S'}$ 的\Term{maximal-torus}{极大环面}就是 $V_{S'}$ 的\Term{maximal-torus}{极大环面}（由 IX 5.2，以及\Term{algebraically-closed-field}{代数闭域}上从\Term{torus}{环面}到\Term{abelian-variety}{阿贝尔簇}的\Term{homomorphism}{同态}平凡这一事实）。因此，上面用 $G$ 定义的\Term{functor}{函子} $Z$ 与用 $V$ 定义的相同。于是可归结到 $G$ \Term{affine}{仿射}。$\mathcal T$ \Term{representable-functor}{可表示}，且在 $k$ 上必“\Term{essentially-free}{本质自由}”（VIII 6.1），故 $\mathcal T_G$ 在 $G$ 上\Term{essentially-free}{本质自由}；$\mathbb T$ 又是 $\mathcal T_G=G_{\mathcal T}$ 的\Term{closed-subscheme}{闭子概形}（例如由 VIII 5.7），故由 VIII 6.4，$Z=\prod_{\mathcal T_G/G}\mathbb T/\mathcal T_G$ 由 $G$ 的\Term{closed-subscheme}{闭子概形}表示。它因而是 $G$ 的\Term{subgroup-scheme}{子群概形}，我们断言它是\Term{multiplicative-type}{乘法型的}：可假定 $k$ \Term{algebraically-closed}{代数闭}，则 $G$ 有\Term{maximal-torus}{极大环面} $T$，由定义 $Z\subset T$；作为\Term{multiplicative-type}{乘法型}群的\Term{algebraic-subgroup}{代数子群}，$Z$ 为\Term{multiplicative-type}{乘法型}（IX.8）。这证明 $Z$ 为 $G$ 的\Term{reductive-center}{约化中心}。它是 $V$ 的最大\Term{central-subgroup}{中心}\Term{multiplicative-type}{乘法型子群}这一点包含在 8.7 中。

b) 必要性显然。设 $Z$ 为 $G$ 的\Term{multiplicative-type}{乘法型子群}，且对每个 $s\in S$，$Z_s$ 为 $G_s$ 的\Term{reductive-center}{约化中心}，来证 $Z$ 为 $G$ 的\Term{reductive-center}{约化中心}。先需证明 $Z$ 包含在 $G$ 的每个\Term{maximal-torus}{极大环面}中（从而这在任意\Term{base-change}{基变换}后仍成立），这是 8.7 的直接推论。再需证明：若 $g$ 是 $G$ 在 $S$ 上的\Term{section}{截面}，且对 $S$ 上任意 $S'$ 及 $G_{S'}$ 的任意\Term{maximal-torus}{极大环面} $T'$，$g_{S'}$ 为 $T'$ 的\Term{section}{截面}，则 $g$ 为 $Z$ 的\Term{section}{截面}。注意，$\prod_{\mathcal T_G/G}\mathbb T/\mathcal T_G=Z'$ 是 fpqc \Term{sheaf}{层}且与环的\Term{inductive-limit}{归纳极限}交换，故可照常先归结到 $S$ \Term{affine}{仿射}，再到 $S$ \Term{noetherian}{诺特}，最后到 $S$ 为\Term{artinian}{阿廷的}。由 7.1 d)，$\mathcal T$ 由 $S$ 上\Term{smooth}{光滑}\Term{prescheme}{预概形}表示；仿照 a)，可见 $\prod_{\mathcal T_G/G}\mathbb T/\mathcal T_G$ 由 $G$ 的\Term{closed-subscheme}{闭子概形} $Z'$ 表示。已经知道 $Z\subset Z'$；又由对 $Z$ 的假设，$Z_k=Z'_k$（$k$ 为\Term{residue-field}{剩余域}），而 $Z$ 在 $S$ 上\Term{flat}{平坦}，故 $Z=Z'$，即 $Z$ 为 $G$ 的\Term{reductive-center}{约化中心}。b) 的其余断言包含在 8.7 中。

c) 立即归结为 d)。
\SourcePage{183}
d) 当然，$\mathfrak g$ 表示 $G$ 的\Term{lie-algebra}{李代数}，$T\to\operatorname{GL}(\mathfrak g)$ 为 $G$ 的\Term{adjoint-representation}{伴随表示}诱导的\Term{homomorphism}{同态}。显然有
\[
T\cap\operatorname{Centr}_G\subset\operatorname{Ker}(T\to\operatorname{GL}(\mathfrak g)).
\]
反向包含的证明与 4.7 d) 相同，此处不再重复。令 $Z$ 为该群；作为 $\operatorname{Ker}(T\to\operatorname{GL}(\mathfrak g))$，它为\Term{multiplicative-type}{乘法型}（例如参见 IX 6.8）。为证它是 $G$ 的\Term{reductive-center}{约化中心}，由 b) 可归结到 $S$ 为域的谱。令 $Z'$ 为\Term{reductive-center}{约化中心}（由 a) 存在），显然 $Z'\subset T\cap\operatorname{Centr}(G)=Z$。另一方面，$Z$ 为包含在\Term{smooth}{光滑}\Term{connected}{连通}\Term{affine}{仿射}\Term{subgroup}{子群} $T$ 中的\Term{central-subgroup}{中心}\Term{multiplicative-type}{乘法型子群}，由 a)，$Z'\supset Z$，故 $Z'=Z$。

e) 在 7.1 f) 的条件下，令 $Z=\operatorname{Ker}u$，并假定对 $S$ 上每个\Term{algebraically-closed-field}{代数闭域} $k$，$Z_k$ 包含在 $G_k$ 的某个\Term{maximal-torus}{极大环面}中。于是容易看出，$T'\mapsto u^{-1}(T')$ 诱导 $G'$ 的\Term{maximal-torus}{极大环面}\Term{set}{集合}与 $G$ 的\Term{maximal-torus}{极大环面}\Term{set}{集合}之间的一一对应。将此用于 8.8 e) 的情形，立即得到所需结论。
\end{proof}
\textbf{注记 8.9。} a) 上述 8.8 的证明不依赖第 4 节的结果，特别地，8.8 a) 的证明未使用 4.4（后者的证明颇费周折）。

b) 容易看出 8.6 中的 $G$ 的\Term{subgroup}{子群} $Z$ 总在\Term{center}{中心}中（无论它是否\Term{representable-functor}{可表示}），因而很自然会想在所有情形下都称它为 $G$ 的\Term{reductive-center}{约化中心}。

c) 也可推广 4.9，得到以下命题：设 $G$ 为\Term{algebraically-closed-field}{代数闭域}上的\Term{smooth}{光滑}\Term{connected}{连通}\Term{algebraic-group}{代数群}，则 $G$ 为\Term{abelian-variety}{阿贝尔簇}被\Term{smooth}{光滑}\Term{connected}{连通}\Term{unipotent}{幺幂}\Term{algebraic-group}{代数群}的\Term{extension}{扩张}（即 $G$ 的\Term{reductive-rank}{约化秩}为零，参见 8.3），当且仅当 $G$ 的\Term{reductive-center}{约化中心}为\Term{unit-group}{单位群}，且 $G$ 的\Term{lie-algebra}{李代数}\Term{nilpotent}{幂零}。

\SGAsection{9}{补充：群概形的作用与不动点}
本节于 2008 年 1 月补入，旨在研究\Term{group-scheme}{群概形}\Term{action}{作用}下的\Term{fixed-point}{不动点}。

\textbf{9.1. \Term{fixed-point}{不动点}\Term{functor}{函子}的\Term{representability}{可表示性}。} 设 $S$ 为\Term{scheme}{概形}，$G$ 为\Term{action}{作用}在 $S$-\Term{scheme}{概形} $X$ 上的 $S$-\Term{group-scheme}{群概形}。定义 $X$ 的\Term{subfunctor}{子函子} $X^G$，即 $X$ 在 $G$ 下的\Term{fixed-point}{不动点函子}：对每个 $S$-\Term{scheme}{概形} $T$，$X^G(T)$ 是 $X(T)$ 中由 $G$ 下的\Term{fixed-point}{不动点}组成的\Term{subset}{子集}（VIII.6.e）。

回顾 [G-R] 第 3.3 节和 [R] 引入的“$S$-\Term{pure}{纯概形}”概念。设 $G$ 为\Term{flat}{平坦}\Term{finite-presentation}{有限表现}的 $S$-\Term{group-scheme}{群概形}。因为 $G$ 的\Term{fiber}{纤维}没有\Term{embedded-component}{嵌入分支}，$G$ 的\Term{purity}{纯性}概念可表述如下：\Term{scheme}{概形} $G$ 是 $S$-\Term{pure}{纯的}，若对每个\Term{closed-point}{闭点}为 $t$ 的\Term{strictly-henselian}{严格 Hensel}\Term{local}{局部} $S$-\Term{scheme}{概形} $T$，$G\times_ST$ 的某一\Term{fiber}{纤维}的每个\Term{generic-point}{泛点} $\widetilde x$ 都\Term{specialization}{特化}到 $G_t$ 的某一点，即 $\widetilde x$ 在 $G\times_ST$ 中的\Term{closure}{闭包}与 $G_t$ 相交。
\SourcePage{184}
\textbf{例。} (1) 若 $G$ 在 $S$ 上\Term{quasi-finite}{拟有限}且\Term{separated}{分离}，则 $G$ 是 $S$-\Term{pure}{纯的}，当且仅当 $G$ 在 $S$ 上\Term{finite}{有限}。

(2) 若 $S=\operatorname{Spec}(R)$ 且 $G=\operatorname{Spec}(A)$，则 $G$ 是 $S$-\Term{pure}{纯的}，当且仅当 $A$ 为\Term{projective-module}{投射} $R$-模（[G-R] 3.3.5）。特别地，若 $G$ \Term{diagonalizable}{可对角化}，则 $G$ 是 $S$-\Term{pure}{纯的}。

(3) \Term{multiplicative-type}{乘法型}\Term{group-scheme}{群概形}是 $S$-\Term{pure}{纯的}，因为\Term{purity}{纯性}概念对于 $S$ 上的\Term{etale-topology}{平展拓扑}是\Term{local-scheme}{局部}的。

(4) 若 $G$ 是 $S$-\Term{proper}{固有的}，或 $G$ 的\Term{fiber}{纤维}\Term{irreducible}{不可约}，或 $S$ 为半\Term{local-scheme}{局部}\Term{artinian}{阿廷的}，则\Term{scheme}{概形} $G$ 是 $S$-\Term{pure}{纯的}。

(5) 特别地，\Term{reductive-group-scheme}{约化} $S$-\Term{group-scheme}{群概形} $G$ 是 $S$-\Term{pure}{纯的}。
\begin{SGAresult}{命题}{9.2}
设 $G$ 为\Term{finite-presentation}{有限表现}、$S$-\Term{flat}{平坦}且 $S$-\Term{pure}{纯}的 $S$-\Term{group-scheme}{群概形}，\Term{action}{作用}在\Term{finite-presentation}{有限表现}的 $S$-\Term{scheme}{概形} $X$ 上。则 $X$ 在 $G$ 下的\Term{fixed-point}{不动点}\Term{functor}{函子} $X^G$ 由 $X$ 的一个在 $S$ 上\Term{finite-presentation}{有限表现}的\Term{closed-subscheme}{闭子概形}表示。
\end{SGAresult}
\begin{proof}
设 $a\in X(S)$，$H$ 是 $G$ 中 $a$ 的\Term{stabilizer}{稳定子}\Term{subgroup-scheme}{子群概形}。则 $a$ 在 $G$ 下不动，当且仅当 $H=G$。而 $H$ 与 $G$ 重合的 $S$ 的\Term{subfunctor}{子函子} $C$（准确地说，对 $S$-\Term{scheme}{概形} $T$，若 $H\times_ST\to G\times_ST$ 为\Term{bijection}{双射}，则 $C(T)=\{\varnothing\}$，否则 $C(T)=\varnothing$）由 $S$ 的一个\Term{closed-subscheme}{闭子概形}表示，且由\Term{finite-type}{有限型}\Term{ideal-sheaf}{理想层}定义（[G-R] 4.1.1）。将这一结果应用于 $S=X$，取 $a$ 为 $X$ 的\Term{generic-point}{泛点}，即 $X(X)=\operatorname{Hom}_S(X,X)$ 中的\Term{identity-map}{恒等映射}即可。
\end{proof}

\textbf{9.3. \Term{infinitesimal}{无穷小}\Term{obstruction}{障碍}。} 假定 $G$ 是\Term{flat}{平坦}\Term{finite-presentation}{有限表现}的 $S$-\Term{group-scheme}{群概形}，（从左）\Term{action}{作用}于 $S$ 上\Term{smooth}{光滑}的\Term{prescheme}{预概形} $X$。下面研究\Term{functor}{函子} $X^G$ 的\Term{formal-smoothness}{形式光滑性}。

这里假定 $S$ 带有\Term{closed-subscheme}{闭子概形} $S_0$，由平方为零的\Term{quasi-coherent}{拟凝聚}\Term{ideal-sheaf}{理想层} $\mathcal I$ 定义。记 $j:S_0\to S$，$G_0=G\times_SS_0$，$X_0=X\times_SS_0$。

设 $\varepsilon_0$ 是 $X^G(S_0)$ 中\Term{liftable}{可提升}为 $X(S)$ 中一点的点。我们将研究将 $\varepsilon_0$ \Term{lifting}{提升}为 $X^G(S)$ 中一点的\Term{obstruction}{障碍}。由于 $X$ \Term{smooth}{光滑}，已知 $\varepsilon_0$ 到 $X(S)$ 的各\Term{lifting}{提升}构成\Term{abelian-group}{阿贝尔群}
\[
\operatorname{Hom}_{\mathcal O_{S_0}}(\varepsilon_0^*(\Omega^1_{X_0/S_0}),\mathcal I)
\]
下的平凡\Term{principal-homogeneous-space}{主齐性空间}（III.0.2、0.3）。于是令
\[
L_0=\mathcal H\!om_{\mathcal O_{S_0}}(\varepsilon_0^*(\Omega^1_{X_0/S_0}),\mathcal I),
\]
这是一个 $\mathcal O_{S_0}$-模。又因 $\varepsilon_0$ 在 $G$ 下不动，$L_0$ 自然为 $G_0$-$\mathcal O_{S_0}$-模。以 $\rho_0:G_0\to\operatorname{Aut}_{S_0}(L_0)$（相应地，$\rho:G\to\operatorname{Aut}_S(j_*L_0)$）表示相应的表示。
\begin{SGAresult}{\Term{lemma}{引理}}{9.4}
存在一个由 $\varepsilon_0$ \Term{canonical}{典范}确定的类 $c(\varepsilon_0)\in\mathrm H^1(G,j_*L_0)\simeq\mathrm H^1(G_0,L_0)$，使 $\varepsilon_0$ \Term{liftable}{可提升}到 $X^G(S)$，当且仅当 $c(\varepsilon_0)=0$。
\end{SGAresult}
\SourcePage{185}
\begin{proof}
\Term{adjunction-isomorphism}{伴随同构} $\mathrm H^1(G,j_*L_0)\simeq\mathrm H^1(G_0,L_0)$ 是\Term{lemma}{引理} III.1.1.2 中的\Term{isomorphism}{同构}。我们在 $S$ 上的小 fppf \Term{site}{景}中工作。对 $S$ 上任意\Term{flat}{平坦}\Term{finite-presentation}{有限表现}\Term{scheme}{概形} $T$，记 $G_T=G\times_ST$、$X_T=X\times_ST$ 为 $T$ 上的相应\Term{object}{对象}。考虑 $S$ 的小 fppf \Term{site}{景}上的\Term{sheaf}{层} $\mathcal A$，使得对每个在 $S$ 上\Term{flat}{平坦}\Term{finite-presentation}{有限表现}的 $T$，有
\[
\mathcal A(T)=\{\varepsilon_0\times_ST\text{ 到 }X(T)\text{ 中的提升}\}.
\]
$L_0$ 的构造与\Term{flat}{平坦}\Term{base-change}{基变换}交换，故前面的讨论表明，$\mathcal A(T)$ 是\Term{abelian-group}{阿贝尔群} $\mathrm H^0(T,j_*L_0)$ 下的平凡\Term{principal-homogeneous-space}{主齐性空间}。仍因 $\varepsilon_0$ 不动，对 $S$ 上每个\Term{flat}{平坦}\Term{finite-presentation}{有限表现}的 $T$ 以及每个 $g\in G(T)$，$g$ 通过\Term{affine}{仿射}\Term{automorphism}{自同构}\Term{action}{作用}于 $\mathcal A(T)$，并与 $G$ 在 $j_*L_0$ 上的\Term{action}{作用}相容。也就是说，
\[
g(a_T+v)=g(a_T)+\rho(g)(v),\qquad v\in\mathrm H^1(T,j_*L_0).
\]
$G$ \Term{flat}{平坦}且\Term{finite-type}{有限型}，故可取 $T=G$、$g$ 为 $G$ 的\Term{generic-point}{泛点} $\operatorname{id}_G$ 来应用上述讨论，从而得到 fppf \Term{sheaf}{层} $G$ 在 $\mathcal A$ 上的\Term{action}{作用}。

设 $a$ 是 $\varepsilon_0$ 在 $X(S)$ 中的一个\Term{lifting}{提升}。记 $g^\sharp$ 为 $G$ 的\Term{generic-point}{泛点}，并以下式定义 $v^\sharp\in\mathrm H^0(G,j_*L_0)$：
\[
\rho(g^\sharp)(v^\sharp)=g^\sharp\cdot a_G-a_G.
\]
对任意 $S$-\Term{prescheme}{预概形} $Y$，令
\[
z(Y):G(Y)\longrightarrow\mathrm H^0(Y,j_*L_0),\qquad g\longmapsto v^\sharp(g).
\]
这定义了 $Z^1(G,L)$ 中的 $1$-\Term{cocycle}{上闭链} $z$（第 I 讲第 5 节）。

它在 $\mathrm H^1(G,j_*L_0)$ 中的类 $c(\varepsilon_0)$ 不依赖于 $a$ 的选择。特别地，若 $\varepsilon_0$ \Term{liftable}{可提升}到 $X^G(S)$，则 $c(\varepsilon_0)=0$。反之，若 $c(\varepsilon_0)=0$，则存在 $w\in\mathrm H^0(S,j_*L_0)$，使对任意 $S$-\Term{prescheme}{预概形} $Y$ 及任意 $g\in G(Y)$，有 $z(Y)(g)=g\cdot w_Y-w_Y$。将此用于 $Y=G$ 和 $g^\sharp$，便得 $a-w\in X^G(S)$。因此已经证明，$\varepsilon_0$ \Term{liftable}{可提升}到 $X^G(S)$，当且仅当 $c(\varepsilon_0)=0$。
\end{proof}
\textbf{注记 9.5。} 当 $S$ \Term{affine}{仿射}，$G$ 在 $S$ 上\Term{affine}{仿射}\Term{flat}{平坦}\Term{finite-type}{有限型}时，可将这一\Term{obstruction}{障碍}细化为类 $\widetilde c(\varepsilon_0)\in\mathrm H^1_{G_0}(X_0,L_0)$，其中 $\mathrm H^1_{G_0}(X_0,L_0)$ 为 Wevers 定义的 $G_0$-\Term{equivariant}{等变}\Term{cohomology}{上同调}群（[W]，附录 C）。这时无需假定 $\varepsilon_0$ \Term{liftable}{可提升}到 $X(S)$。

\textbf{9.6. \Term{fixed-point}{不动点}的\Term{smoothness}{光滑性}。}
\begin{SGAresult}{\Term{theorem}{定理}}{9.7}
设 $G$ 是\Term{noetherian}{诺特}\Term{base}{基底} $S$ 上的\Term{flat}{平坦}\Term{finite-type}{有限型} $S$-\Term{group-scheme}{群概形}，\Term{action}{作用}在 $S$ 上\Term{smooth}{光滑}的\Term{scheme}{概形} $X$ 上。假定 $G$ 是 $S$-\Term{pure}{纯的}，且对 $S$ 上每个\Term{geometric-point}{几何点} $s$（其\Term{residue-field}{剩余域}\Term{algebraically-closed}{代数闭}），有
\[
\mathrm H^1\bigl(G_s,(\Omega^1_{X_s})^*\bigr)=0.
\]
则 $X$ 在 $G$ 下的\Term{fixed-point}{不动点}\Term{functor}{函子} $X^G$ 由 $X$ 的一个在 $S$ 上\Term{smooth}{光滑}的\Term{closed-subscheme}{闭子概形}表示。
\end{SGAresult}
\begin{proof}
来证明\Term{theorem}{定理} 9.7。已知 $X^G$ 由 $X$ 的一个在 $S$ 上\Term{finite-presentation}{有限表现}的\Term{closed-subscheme}{闭子概形}表示。可假定 $S=\operatorname{Spec}(A)$ \Term{local}{局部}。根据\Term{smoothness-criterion}{光滑性判据}（SGA1 III.3.1.iii.bis），只需对 $S'=\operatorname{Spec}(A')$ 及\Term{closed-subscheme}{闭子概形}
\SourcePage{186}
$S'_0=\operatorname{Spec}(A'/\mathcal I')=\operatorname{Spec}(A'_0)$ 验证 $X^G$ 的\Term{formal-smoothness}{形式光滑性}，其中 $A'$ 是\Term{local}{局部}\Term{artinian}{阿廷}环，$\mathcal I'$ 为 $A'$ 的平方为零的\Term{ideal}{理想}。经逐步剥离，可归结到 $\mathcal I'$ 被 $A'$ 的\Term{maximal-ideal}{极大理想}\Term{annihilate}{零化}的情形。特别地，若 $k$ 表示 $A'$ 的\Term{residue-field}{剩余域}，则 $\mathcal I'$ 为 $k$-\Term{vector-space}{向量空间}。

于是给定 $\varepsilon'_0\in X^G(S'_0)$，希望将它\Term{lifting}{提升}到 $X^G(S')$。记 $X'=X\times_SS'$、$G'=G\times_SS'$、$X'_0=X\times_SS'_0$、$G'_0=G\times_SS'_0$。$S'$ \Term{affine}{仿射}且 $X'$ 在 $S'$ 上\Term{smooth}{光滑}，故 $\varepsilon'_0\in X^G(S'_0)$ \Term{liftable}{可提升}到 $X(S')$。由\Term{lemma}{引理} 9.4，类
\[
c(\varepsilon_0)\in\mathrm H^1(G'_0,L'_0)
\]
是将 $\varepsilon'_0$ \Term{lifting}{提升}到 $X^G(S')$ 的\Term{obstruction}{障碍}，其中
\[
L'_0=\mathcal H\!om_{\mathcal O_{S'_0}}\bigl(\varepsilon'_0{}^*(\Omega^1_{X'_0/S'_0}),\mathcal I'\bigr).
\]
$L'_0$ 为 $k$-\Term{vector-space}{向量空间}，故有\Term{canonical-isomorphism}{典范同构}
\[
\mathrm H^1(G'_0,L'_0)\xrightarrow{\sim}\mathrm H^1(G'_0\times_{A'_0}k,L'_0).
\]
只需证明 $\mathrm H^1(G'_0\times_{A'_0}k,L'_0)=0$。记 $\overline k$ 为 $k$ 的\Term{algebraic-closure}{代数闭包}，有\Term{isomorphism}{同构}（\Term{affine}{仿射}情形见 IX.3.1，一般情形见\Term{lemma}{引理} 9.11）
\[
\mathrm H^1(G'_0\times_{A'_0}k,L'_0)\otimes_k\overline k
\xrightarrow{\sim}\mathrm H^1(G'_0\times_{A'_0}\overline k,L'_0\otimes_k\overline k).
\]
注意表示 $L'_0\otimes_k\overline k$ 是若干个 $(\Omega^1_{X'_0\times_k\overline k})^*$ 的\Term{direct-sum}{直和}。根据假设，$\mathrm H^1(G'_0\times_kA'_0,(\Omega^1_{X'_0\times_k\overline k})^*)=0$，这就\Term{pushout}{推出} $\mathrm H^1(G'_0\times_{A'_0}k,L'_0)=0$。
\end{proof}
\begin{SGAresult}{推论}{9.8}
设 $G$ 是\Term{flat}{平坦}\Term{finite-presentation}{有限表现}的 $S$-\Term{group-scheme}{群概形}，\Term{action}{作用}在 $S$ 上\Term{smooth}{光滑}\Term{finite-presentation}{有限表现}的\Term{scheme}{概形} $X$ 上。假定 $G$ 有一个\Term{composition-series}{合成列}，其因子属于下述类型之一：

(1) $S$-\Term{abelian-scheme}{阿贝尔概形}（即 $G$ 在 $S$ 上\Term{smooth}{光滑}，其\Term{fiber}{纤维}为\Term{abelian-variety}{阿贝尔簇}）；

(2) \Term{multiplicative-type}{乘法型} $S$-\Term{group-scheme}{群概形}；

(3) \Term{finite}{有限}、\Term{etale}{平展}、\Term{degree}{次数}在 $S$ 上\Term{invertible}{可逆}的 $S$-\Term{group-scheme}{群概形}；

(4) $S$ 为 $\mathbb Q$-\Term{scheme}{概形}时的 $S$-\Term{reductive-group-scheme}{约化群}。

则 $X$ 在 $G$ 下的\Term{fixed-point}{不动点}\Term{functor}{函子} $X^G$ 由 $X$ 的一个在 $S$ 上\Term{finite-presentation}{有限表现}且\Term{smooth}{光滑}的\Term{closed-subscheme}{闭子概形}表示。
\end{SGAresult}
\begin{proof}
事实上，若 $1\to G'\to G\to G''\to1$ 是满足\Term{theorem}{定理}假设的 $S$-群\Term{exact-sequence}{正合列}，则 $S$-群 $G''$ \Term{action}{作用}于\Term{functor}{函子} $X^{G'}$，而\Term{functor}{函子} $X^G$ 正是 $X^{G'}$ 对 $G''$ 的\Term{fixed-point}{不动点函子}。因此可归结到上述四种类型之一的 $S$-群。通过取\Term{limit}{极限}，可假定 $S$ \Term{noetherian}{诺特}。我们对 $S$ 的某个\Term{geometric-point}{几何点} $s$ 上的\Term{fiber}{纤维} $G_s$ 及 $G_s$ 的一个\Term{finite-dimensional}{有限维}\Term{linear-representation}{线性表示} $V$，验证上述\Term{cohomology}{上同调}判据。

在\Term{abelian-scheme}{阿贝尔概形}情形，从 $G_s^n$ 到 $V$ 的每个映射均为常值，因此对 $i>0$，$\mathrm H^i(G_s,V)$ 均为零。

在\Term{multiplicative-type}{乘法型}\Term{group-scheme}{群概形}情形，$G_s$ 为\Term{diagonalizable}{可对角化}群。\Term{theorem}{定理} I.5.3.3 表明，对 $i>0$，$\mathrm H^i(G_s,V)$ 均为零。

后两种情形的论证相同，都基于 $G_s$ 的\Term{linear-representation}{线性表示}的\Term{semisimplicity}{半单性}。事实上，对于\Term{finite}{有限}、\Term{etale}{平展}、\Term{degree}{次数}在 $S$ 上\Term{invertible}{可逆}的 $S$-\Term{group-scheme}{群概形}，
\SourcePage{187}
$G_s$ 为\Term{degree}{次数}在\Term{residue-field}{剩余域} $\kappa(s)$ 中\Term{invertible}{可逆}的\Term{finite}{有限}\Term{constant-group}{常群}；众所周知，$G_s$ 的每个\Term{linear-representation}{线性表示}都\Term{semisimple}{半单}。

在\Term{reductive-group-scheme}{约化}情形，假定（\Term{algebraically-closed}{代数闭}的）\Term{residue-field}{剩余域} $\kappa(s)$ 的\Term{characteristic}{特征}为零。这时已知 $G_s$ 的每个\Term{linear-representation}{线性表示}都\Term{semisimple}{半单}（见 [T-Y]，\Term{theorem}{定理} 27.3.3）。
\end{proof}
\textbf{注记 9.9。} 这特别适用于 $G/S$ \Term{action}{作用}在 $S$ 上\Term{smooth}{光滑}\Term{group-scheme}{群概形} $H$ 上的情形。若 $G$ 为\Term{multiplicative-type}{乘法型}群，通过内变换\Term{action}{作用}于 $S$ 上\Term{smooth}{光滑}且\Term{affine}{仿射}的 $H$，便重新得到推论 XI.5.3，即 $G$ 在 $H$ 中的\Term{centralizer}{中心化子}\Term{smooth}{光滑}。

\textbf{注记 9.10。} 当 $S=\operatorname{Spec}(k)$，$G$ 为\Term{algebraically-closed-field}{代数闭域} $k$ 上的\Term{linearly-reductive}{线性约化}\Term{algebraic-group}{代数群}时，这一结果由 Fogarty（[F]，\Term{theorem}{定理} 5.4）和 Iversen（[I]，命题 1.3）分别独立得到。

\Term{theorem}{定理} 9.7 的证明使用了下述\Term{lemma}{引理}；在\Term{affine}{仿射}\Term{group-scheme}{群概形}的情形，它是熟知的（IX.3.1）。
\begin{SGAresult}{\Term{lemma}{引理}}{9.11}
设 $G$ 是\Term{affine}{仿射}\Term{scheme}{概形} $S=\operatorname{Spec}(A)$ 上的\Term{group-scheme}{群概形}，$f:S'=\operatorname{Spec}(A')\to S$ 为\Term{flat}{平坦}\Term{morphism}{态射}。令 $G'=G\times_SS'$。则对每个\Term{quasi-coherent}{拟凝聚} $G$-$\mathcal O_S$-模 $M$，有\Term{canonical-isomorphism}{典范同构}
\[
\mathrm H^i(G,M)\otimes_AA'\xrightarrow{\sim}\mathrm H^i(G',M\otimes_AA')\qquad(i\geq0).
\]
\end{SGAresult}
\begin{proof}
Hochschild \Term{cohomology}{上同调}的 $n\geq1$ 次链由它们在点 $(\operatorname{id},\ldots,\operatorname{id})\in G(G)\times\cdots\times G(G)$ 处的值确定，故有 $A$-模\Term{isomorphism}{同构}
\[
C^n(G,M)\simeq\Gamma(G^n,M\otimes_A\mathcal O_{G^n}).
\]
\Term{cohomology}{上同调}群 $\mathrm H^n(G,M)$ 为\Term{complex}{复形}
\[
L\longrightarrow\Gamma(G,M\otimes_A\mathcal O_G)
\longrightarrow\Gamma(G^2,M\otimes_A\mathcal O_{G^2})
\longrightarrow\Gamma(G^3,M\otimes_A\mathcal O_{G^3})\longrightarrow\cdots
\]
的第 $n$ 个\Term{cohomology}{上同调}群。由于 $A'$ 在 $A$ 上\Term{flat}{平坦}，由 EGA IV$_1$ 1.7.21，有
\[
\Gamma(G^n,M\otimes_A\mathcal O_{G^n})\otimes_AA'
\xrightarrow{\sim}\Gamma\bigl(G'{}^n,(M\otimes_AA')\otimes_{A'}\mathcal O_{G'{}^n}\bigr).
\]
取\Term{cohomology}{上同调}便得所需断言。
\end{proof}

\textbf{参考文献}\footnote{编者注：本讲引用的补充参考文献。}

\noindent [Fo73] J. Fogarty，\textit{\Term{fixed-point-scheme}{不动点概形}}（Fixed point schemes），Amer. J. Math. \textbf{95}（1973），35--51。

\noindent [RG71] M. Raynaud，L. Gruson，\textit{\Term{flatness}{平坦性}与\Term{projective-module}{投射}性的判据}（Critères de platitude et de projectivité），Invent. math. \textbf{13}（1971），1--89。

\noindent [Iv72] B. Iversen，\textit{\Term{torus}{环面}在\Term{algebraic-variety}{代数簇}上\Term{action}{作用}的\Term{fixed-point}{不动点}公式}（A fixed point formula for action of tori on algebraic varieties），Invent. math. \textbf{16}（1972），229--236。
\SourcePage{188}
\noindent [Ray72] M. Raynaud，\textit{\Term{algebraic-geometry}{代数几何}中的\Term{flat-module}{平坦模}}（Flat modules in algebraic geometry），Compositio Math. \textbf{24}（1972），11--31。

\noindent [TY06] P. Tauvel，R. W. T. Yu，\textit{\Term{lie-algebra}{李代数}与\Term{algebraic-group}{代数群}}（Lie algebras and algebraic groups），Springer-Verlag，2006。

\noindent [We05] S. Wewers，\textit{带有\Term{group-scheme}{群概形}\Term{action}{作用}的\Term{curve}{曲线}的\Term{form}{形式}\Term{deformation}{形变}}（Formal deformation of curves with group scheme action），Ann. Inst. Fourier (Grenoble) \textbf{55}（2005），1105--1165。
